%%%%%%%%%%%%%%%%%%%%%%%%%%%%%%%%%%%%%%%%%%%%%%%%%%%%%%%%%%%%%%%%%%%%%%%%%%%%%%%
%% TKS FORMAL MATHEMATICAL MANUAL v7.0 — ACADEMIC EXPANSION
%% Full Integration Release
%% Complete Noetic Fractal Calculus, RPM Monad, Dynamic Semantics,
%% Concurrency Model, and Quantum Noetic Algebra
%%%%%%%%%%%%%%%%%%%%%%%%%%%%%%%%%%%%%%%%%%%%%%%%%%%%%%%%%%%%%%%%%%%%%%%%%%%%%%%

\documentclass[11pt,twoside,openright]{book}

%% ============================================================================
%% PACKAGES (Preserved from v6.1 + v7.0 additions)
%% ============================================================================
\usepackage[utf8]{inputenc}
\usepackage[T1]{fontenc}
\usepackage{amsmath,amssymb,amsthm,amsfonts}
\usepackage{mathtools}
\usepackage{stmaryrd}           % For semantic brackets
\usepackage{bm}                 % Bold math
\usepackage{enumitem}
\usepackage{booktabs}
\usepackage{array}
\usepackage{longtable}
\usepackage{multirow}
\usepackage{graphicx}
\usepackage{tikz}
\usetikzlibrary{cd,arrows,positioning,calc}
\usepackage{hyperref}
\usepackage{cleveref}
\usepackage{xcolor}
\usepackage{tcolorbox}
\tcbuselibrary{theorems,skins,breakable}
\usepackage{fancyhdr}
\usepackage{titlesec}
\usepackage{geometry}
\usepackage{listings}
\usepackage{multicol}

%% v7.0 PACKAGE ADDITIONS
\usepackage{braket}             % For quantum notation
\usepackage{tikz-cd}            % For commutative diagrams
\usepackage{bussproofs}         % For proof trees

\geometry{
  paper=letterpaper,
  inner=1.2in,
  outer=1in,
  top=1in,
  bottom=1in
}

%% ============================================================================
%% THEOREM ENVIRONMENTS (Preserved from v6.1)
%% ============================================================================
\theoremstyle{definition}
\newtheorem{definition}{Definition}[chapter]
\newtheorem{axiom}[definition]{Axiom}

\theoremstyle{plain}
\newtheorem{theorem}[definition]{Theorem}
\newtheorem{lemma}[definition]{Lemma}
\newtheorem{proposition}[definition]{Proposition}
\newtheorem{corollary}[definition]{Corollary}

\theoremstyle{remark}
\newtheorem{remark}[definition]{Remark}
\newtheorem{example}[definition]{Example}
\newtheorem{notation}[definition]{Notation}

%% ============================================================================
%% CUSTOM COMMANDS (Preserved from v6.1)
%% ============================================================================
% Semantic brackets
\newcommand{\sem}[1]{\llbracket #1 \rrbracket}

% TKS-specific symbols
\newcommand{\Worlds}{\mathcal{W}}
\newcommand{\Ideas}{\mathcal{I}}
\newcommand{\Elements}{\mathcal{E}}
\newcommand{\Noetics}{\mathcal{N}}
\newcommand{\Foundations}{\mathcal{F}}
\newcommand{\Acquisitions}{\mathcal{A}}
\newcommand{\Noetica}{\mathbf{Noetica}}
\newcommand{\Acq}{\mathbf{Acq}}

% Noetic operators
\newcommand{\noe}[1]{\nu_{#1}}

% Tootra operations
\newcommand{\Tadd}{\oplus_T}
\newcommand{\Tsub}{\ominus_T}
\newcommand{\Tmul}{\otimes_T}
\newcommand{\Tdiv}{\oslash_T}

% World association
\newcommand{\Wadd}{\oplus_W}
\newcommand{\Wsub}{\ominus_W}

% Type judgments
\newcommand{\types}{\vdash}
\newcommand{\typmark}{:}

% RPM Monad
\newcommand{\RPM}{\mathfrak{R}}
\newcommand{\rpmret}{\mathsf{return}}
\newcommand{\rpmbind}{\mathbin{\gg\!=}}

% Fractal notation
\newcommand{\Frac}{\Phi}
\newcommand{\FracSet}{\mathbf{Frac}}

% Category theory
\newcommand{\Ob}{\mathrm{Ob}}
\newcommand{\Hom}{\mathrm{Hom}}
\newcommand{\id}{\mathrm{id}}
\newcommand{\dom}{\mathrm{dom}}
\newcommand{\cod}{\mathrm{cod}}

% Ordinal function
\newcommand{\ord}{\mathrm{ord}}

% Evaluation relations
\newcommand{\evalto}{\Downarrow}
\newcommand{\smallstep}{\to}
\newcommand{\bigstep}{\Downarrow}

% Acquisition state
\newcommand{\AcqState}{\sigma}
\newcommand{\Satisfied}{\mathsf{Sat}}
\newcommand{\Failed}{\mathsf{Fail}}

%% ============================================================================
%% v7.0 CUSTOM COMMAND ADDITIONS
%% ============================================================================

% Quantum notation
\newcommand{\Hspace}{\mathcal{H}}
\newcommand{\qket}[1]{|#1\rangle}
\newcommand{\qbra}[1]{\langle#1|}
\newcommand{\qbraket}[2]{\langle#1|#2\rangle}
\newcommand{\qop}[1]{\hat{#1}}

% Parallel composition
\newcommand{\pcomp}{\parallel}

% Effect typing
\newcommand{\efftype}[2]{\mathsf{RPM}[#1, #2]}

% Kleisli category
\newcommand{\Kleisli}{\mathcal{K}}

% Subtyping
\newcommand{\subtype}{<:}

% Eigenvalue spectrum
\newcommand{\Spec}{\mathrm{Spec}}

% v7.0 markers for new content
\newcommand{\vnew}{\textcolor{blue!70!black}{\textbf{[NEW v7.0]}}}
\newcommand{\vexp}{\textcolor{green!50!black}{\textbf{[EXPANDED v7.0]}}}

%% ============================================================================
%% HEADER/FOOTER
%% ============================================================================
\pagestyle{fancy}
\fancyhf{}
\fancyhead[LE]{\leftmark}
\fancyhead[RO]{\rightmark}
\fancyfoot[C]{\thepage}
\renewcommand{\headrulewidth}{0.4pt}

%% ============================================================================
%% TITLE INFO
%% ============================================================================
\title{%
  \Huge\bfseries TKS FORMAL MATHEMATICAL MANUAL\\[0.5em]
  \Large Version 7.0 — Academic Expansion\\[0.5em]
  \large\textit{Full Integration Release}\\[1em]
  \normalsize The Complete Rigorous Formalization of the\\
  Tootra Kabbalistic System\\[0.5em]
  with Complete Noetic Fractal Calculus, RPM Monad Theory,\\
  Dynamic Semantics, Concurrency Model, and Quantum Noetic Algebra
}
\author{%
  Compiled from Canonical Sources\\[0.5em]
  \textit{A Graduate-Level Treatise in Formal Metaphysical Mathematics}
}
\date{2024}

%% ============================================================================
%% BEGIN DOCUMENT
%% ============================================================================
\begin{document}

\frontmatter
\maketitle

%% ============================================================================
%% WHAT'S NEW IN v7.0 (NEW SECTION)
%% ============================================================================
\chapter*{What's New in Version 7.0}
\addcontentsline{toc}{chapter}{What's New in Version 7.0}

\begin{tcolorbox}[colback=blue!5!white,colframe=blue!75!black,title=Version 7.0 — Academic Expansion Release]
Version 7.0 represents a major expansion of the TKS formal system, building upon the solid foundation of v6.1 (Dynamic Semantics Release) with significant new theoretical frameworks and complete formalization of all major subsystems.
\end{tcolorbox}

\section*{Major Additions}

\subsection*{1. Complete Noetic Fractal Calculus Integration}

\begin{itemize}
  \item \textbf{Fractal Type Theory:} Formal type constructors $\mathsf{Frac}[\langle k_1:\cdots:k_n \rangle]$ with complete kinding rules and type-level sequence operations.
  \item \textbf{Eigenvalue Analysis:} Spectrum analysis for fractals, fixed point theorems, and stability classification with rigorous proofs.
  \item \textbf{Convergence Theorems:} Complete proofs of fractal convergence under iteration, with rate bounds and basin of attraction characterization.
  \item \textbf{Fractal Categories:} Category-theoretic treatment of fractal spaces with functors between them.
\end{itemize}

\subsection*{2. RPM as a Complete Monad System}

\begin{itemize}
  \item \textbf{Kleisli Category:} Full formalization of $\Kleisli_\RPM$ with explicit morphism composition.
  \item \textbf{Effect System:} RPM formalized as an algebraic effect with handlers and effect polymorphism.
  \item \textbf{Multi-Goal Structures:} Parallel prerequisite checking with synchronization semantics.
  \item \textbf{Natural Transformations:} Relationships between RPM and other computational monads.
\end{itemize}

\subsection*{3. Comprehensive Dynamic Semantics}

\begin{itemize}
  \item \textbf{Complete Evaluation Rules:} Big-step ($\bigstep$) and small-step ($\smallstep$) semantics for \emph{all} TKS constructs including Elements, Noetics, Tootra operations, Fractals, ACBE, and RPM.
  \item \textbf{Evaluation Environments:} Formal treatment of variable bindings, stores, acquisition states, and fractal contexts.
  \item \textbf{State Threading:} Explicit state monad formalization for acquisition tracking.
  \item \textbf{Confluence and Determinism:} Proofs that evaluation is deterministic for well-typed programs.
\end{itemize}

\subsection*{4. Expanded Type System}

\begin{itemize}
  \item \textbf{Complete Type Constructors:} Full formalization of $\mathsf{Domain}$, $\mathsf{Aspect}$, $\mathsf{Foundation}$, $\mathsf{SubFoundation}$, $\mathsf{Noe}[k]$, $\mathsf{Frac}[\bar{k}]$, and $\mathsf{RPM}[T]$.
  \item \textbf{Full Inference Rules:} Complete set of typing rules for all language constructs.
  \item \textbf{Progress and Preservation:} Complete proofs (not just statements) of type safety.
  \item \textbf{Polymorphism:} System F-style parametric polymorphism with type abstraction and application.
  \item \textbf{Subtyping Lattice:} Complete lattice structure with reflexivity, transitivity, and antisymmetry proofs.
\end{itemize}

\subsection*{5. Denotational Semantics Revision}

\begin{itemize}
  \item \textbf{Extended Semantic Function:} $\sem{\cdot}$ defined for all fractal and monadic constructs.
  \item \textbf{Semantic Domains:} Complete domain theory for TKS values including lifted domains and continuous functions.
  \item \textbf{Soundness Theorem:} Full proof that operational semantics agrees with denotational semantics.
  \item \textbf{Adequacy:} Proof that denotational equality implies observational equivalence.
\end{itemize}

\subsection*{6. Concurrency Model (Prototype)}

\begin{itemize}
  \item \textbf{Parallel Fractal Evaluation:} Formal semantics for concurrent fractal application.
  \item \textbf{Synchronization Primitives:} Barriers, locks, and channels for coordinating fractal threads.
  \item \textbf{Determinacy Analysis:} Conditions under which parallel evaluation is deterministic.
  \item \textbf{Race Condition Detection:} Static analysis for detecting potential conflicts.
\end{itemize}

\subsection*{7. Quantum Noetic Algebra (Prototype)}

\begin{itemize}
  \item \textbf{Hilbert Space Formulation:} Noetics as linear operators on quantum Idea space $\Hspace_\Ideas$.
  \item \textbf{Superposition:} Quantum superposition of Ideas with probability amplitudes.
  \item \textbf{Measurement Semantics:} Projection postulate applied to Noetic observation.
  \item \textbf{Commutation Relations:} Uncertainty principles for non-commuting Noetics.
  \item \textbf{Quantum Fractals:} Unitary operators as quantum fractal evolutions.
\end{itemize}

\subsection*{8. Comprehensive Worked Examples}

\begin{itemize}
  \item \textbf{20+ Complete Examples:} Detailed evaluation traces for diverse scenarios.
  \item \textbf{RPM Resolution Steps:} Step-by-step prerequisite checking with state evolution.
  \item \textbf{Fractal Application Sequences:} Multi-step fractal evaluation with intermediate results.
  \item \textbf{Type Derivation Trees:} Complete typing derivations for complex expressions.
\end{itemize}

\section*{Relationship to Prior Versions}

\begin{center}
\begin{tabular}{lll}
\toprule
\textbf{Version} & \textbf{Release Name} & \textbf{Key Features} \\
\midrule
v3.2.1 & Initial Formalization & Core definitions \\
v5.0 & Mathematical Framework & Category theory, type theory \\
v6.0 & Academic Release & Denotational semantics, proofs \\
v6.1 & Dynamic Semantics & RPM monad, fractal calculus (partial) \\
\textbf{v7.0} & \textbf{Academic Expansion} & \textbf{Complete integration of all systems} \\
\bottomrule
\end{tabular}
\end{center}

\textbf{Compatibility:} v7.0 is a strict superset of v6.1. All v6.1 definitions, theorems, and proofs are preserved unchanged. v7.0 adds new material without modifying existing content.

%% ============================================================================
%% ABSTRACT
%% ============================================================================
\chapter*{Abstract}
\addcontentsline{toc}{chapter}{Abstract}

This manual presents the complete formal mathematical specification of the \textbf{Tootra Kabbalistic System (TKS)}, Version 7.0---the \textbf{Academic Expansion Release}. Building upon the canonical v6.1 Dynamic Semantics Release, this edition provides the most comprehensive formalization of TKS to date:

\begin{enumerate}[label=(\arabic*)]
  \item A \textbf{complete ontology} of 40 Elements across 4 Worlds with 10 Noetic operators (preserved from v6.0/v6.1).

  \item A \textbf{formal algebra} with Noetic composition, Foundation filtering, and Tootra arithmetic (preserved from v6.0/v6.1).

  \item A \textbf{category-theoretic foundation} establishing TKS as the category $\Noetica$ (preserved from v6.0/v6.1).

  \item \vnew\ A \textbf{complete Noetic Fractal Calculus} with:
    \begin{itemize}
      \item Fractal type theory with indexed types $\mathsf{Frac}[\langle k_1:\cdots:k_n \rangle]$
      \item Eigenvalue spectrum analysis and fixed point theorems
      \item Convergence proofs with rate bounds
      \item Fractal categories and functorial relationships
    \end{itemize}

  \item \vnew\ A \textbf{complete RPM Monad theory} with:
    \begin{itemize}
      \item Kleisli category $\Kleisli_\RPM$ formalization
      \item Algebraic effect system interpretation
      \item Multi-goal parallel structures
      \item Natural transformations to other monads
    \end{itemize}

  \item \vnew\ \textbf{Comprehensive dynamic semantics} including:
    \begin{itemize}
      \item Big-step and small-step rules for all constructs
      \item Formal evaluation environments
      \item State threading with explicit monad structure
      \item Confluence and determinism proofs
    \end{itemize}

  \item \vnew\ An \textbf{expanded type system} with:
    \begin{itemize}
      \item Complete type constructors for all TKS entities
      \item Full Progress and Preservation proofs
      \item System F-style polymorphism
      \item Subtyping lattice with soundness proof
    \end{itemize}

  \item \vnew\ A \textbf{Concurrency Model prototype} for parallel fractal evaluation.

  \item \vnew\ A \textbf{Quantum Noetic Algebra prototype} treating Noetics as operators on Hilbert space.

  \item \textbf{20+ fully worked examples} with complete evaluation traces, type derivations, and RPM resolution steps.

  \item \textbf{Rigorous proofs} of all major theorems with academic-level detail.
\end{enumerate}

Version 7.0 represents the \textbf{Academic Expansion Release}, providing:
\begin{itemize}
  \item Complete formalization suitable for proof assistant implementation (Coq, Agda, Lean)
  \item Sufficient detail for peer-reviewed academic publication
  \item Comprehensive reference for TKS interpreter/compiler implementation
  \item Foundation for future extensions including distributed and quantum TKS
\end{itemize}

\vspace{1em}
\noindent\textbf{Keywords:} Formal metaphysics, category theory, type theory, denotational semantics, operational semantics, monads, Kleisli categories, algebraic effects, fractal calculus, eigenfractals, quantum algebra, Hilbert spaces, concurrency, Kabbalistic mathematics, consciousness modeling, transformation algebra.

%% ============================================================================
%% TABLE OF CONTENTS
%% ============================================================================
\tableofcontents

%% ============================================================================
%% PREFACE
%% ============================================================================
\chapter*{Preface}
\addcontentsline{toc}{chapter}{Preface}

The Tootra Kabbalistic System (TKS) represents an ambitious attempt to formalize metaphysical concepts using the rigorous tools of modern mathematics. This manual serves as the definitive reference for TKS v7.0, the \textbf{Academic Expansion Release}.

\section*{Goals of Version 7.0}

Version 7.0 aims to provide:

\begin{enumerate}
  \item \textbf{Completeness:} Every construct in TKS has formal semantics---no informal gaps remain.

  \item \textbf{Rigor:} All theorems have complete proofs suitable for formalization in proof assistants.

  \item \textbf{Extensibility:} The framework supports future extensions including concurrency and quantum interpretations.

  \item \textbf{Implementability:} Sufficient specification detail for correct interpreter/compiler construction.
\end{enumerate}

\section*{What's Preserved from v6.1}

All content from v6.1 is preserved unchanged:
\begin{itemize}
  \item Canonical ontology (Worlds, Elements, Noetics, Foundations, Acquisitions)
  \item Algebraic structures (Noetic algebra, Tootra arithmetic)
  \item Category Noetica and ACBE functor
  \item Basic RPM monad structure and laws
  \item Basic fractal definitions and simplification rules
  \item Type system foundations
  \item Language specification (BNF grammar)
\end{itemize}

\section*{What's New in v7.0}

Major additions include:
\begin{itemize}
  \item Complete fractal type theory with eigenfractal analysis
  \item RPM Kleisli category and effect system formalization
  \item Full dynamic semantics for all constructs
  \item Complete Progress and Preservation proofs
  \item Concurrency model prototype
  \item Quantum Noetic algebra prototype
  \item Doubled number of worked examples
\end{itemize}

\section*{Intended Audience}

This manual is written for:
\begin{itemize}
  \item Graduate students in mathematics, computer science, or philosophy
  \item Researchers interested in formal approaches to metaphysics
  \item Implementers building TKS-based software systems
  \item Proof assistant users formalizing TKS in Coq/Agda/Lean/Isabelle
  \item Scholars studying the intersection of Kabbalistic thought and formal methods
\end{itemize}

\section*{Prerequisites}

Readers should have familiarity with:
\begin{itemize}
  \item Basic set theory and abstract algebra
  \item Category theory (categories, functors, natural transformations, monads, Kleisli categories)
  \item Type theory fundamentals (simply-typed lambda calculus, System F)
  \item Programming language semantics (denotational and operational)
  \item Basic linear algebra (for quantum chapter)
\end{itemize}

\section*{Document Structure}

The manual is organized into nine major parts:
\begin{enumerate}
  \item \textbf{Foundations}: Ontology, Elements, Noetics, Foundations
  \item \textbf{Algebraic Structures}: Noetic algebra, Tootra arithmetic
  \item \textbf{Categorical Framework}: Category Noetica, ACBE functor
  \item \textbf{RPM and Fractal Calculus}: Complete monad theory and fractal type theory
  \item \textbf{Complete Dynamic Semantics}: Evaluation rules for all constructs
  \item \textbf{Type Theory}: Complete type system with proofs
  \item \textbf{Advanced Topics}: Concurrency and quantum prototypes
  \item \textbf{Language and Implementation}: Specification and architecture
  \item \textbf{Applications}: Comprehensive examples and reference
\end{enumerate}

\mainmatter

%% ============================================================================
%% PART I: FOUNDATIONS
%% [CONTENT PRESERVED FROM v6.1 - TO BE COPIED]
%% ============================================================================
\part{Foundations}

\chapter{Introduction}
\label{ch:introduction}

\section{Aims and Scope of TKS}

The Tootra Kabbalistic System (TKS) is a formal mathematical framework designed to model the structure of consciousness, transformation, and manifestation. Unlike purely empirical theories, TKS operates in the domain of \emph{formal metaphysics}---it provides a rigorous algebraic and categorical language for reasoning about concepts traditionally relegated to philosophy or mysticism.

\begin{definition}[Formal Metaphysics]
\label{def:formal-metaphysics}
A \textbf{formal metaphysical system} is a mathematical structure $\mathcal{S} = (\mathcal{O}, \mathcal{M}, \mathcal{A})$ where:
\begin{itemize}
  \item $\mathcal{O}$ is a set of \emph{ontological primitives} (the basic entities of the system),
  \item $\mathcal{M}$ is a set of \emph{morphisms} or transformations between entities,
  \item $\mathcal{A}$ is a set of \emph{axioms} governing the behavior of primitives and morphisms.
\end{itemize}
\end{definition}

TKS instantiates this schema with:
\begin{itemize}
  \item $\mathcal{O}$: The 40 Elements, 4 Worlds, 7 Foundations, and the Idea space $\Ideas$.
  \item $\mathcal{M}$: The 10 Noetic operators $\noe{0}, \ldots, \noe{9}$, the RPM monad $\RPM$, and Noetic fractals.
  \item $\mathcal{A}$: The axioms of composition, duality, precedence, prerequisite satisfaction, and fractal convergence.
\end{itemize}

\section{The Central Thesis}

TKS is built on a single foundational claim:

\begin{tcolorbox}[colback=blue!5!white,colframe=blue!75!black,title=Central Thesis of TKS]
\textbf{All phenomena of consciousness, transformation, and manifestation can be modeled as algebraic operations on a structured space of Ideas, mediated by a finite set of Noetic operators, organized across a hierarchy of Worlds, filtered through life-domain Foundations, subject to prerequisite constraints (RPM), and capable of recursive self-application (Fractals).}
\end{tcolorbox}

This thesis is \emph{not} an empirical claim about physics or neuroscience. Rather, it is a \emph{formal postulate} that enables systematic reasoning about transformation processes.

\section{v6.1 Dynamic Semantics: Overview}

Version 6.1 introduces two major components that provide TKS with full \emph{dynamic semantics}:

\subsection{The RPM Monad}

The Recursive Prerequisite Model (RPM) is formalized as a monad $\RPM$ on the category of TKS expressions with acquisition state. This provides:
\begin{itemize}
  \item A formal mechanism for sequencing computations that depend on prerequisite satisfaction.
  \item Failure propagation: if a prerequisite fails, the entire computation fails gracefully.
  \item A clean separation between pure TKS expressions and effectful prerequisite checking.
\end{itemize}

\subsection{Noetic Fractal Calculus}

Noetic Fractals are formalized as higher-order operators with:
\begin{itemize}
  \item A well-defined type system.
  \item Composition and simplification rules.
  \item Convergence criteria for iterated application.
  \item Classification into stabilizing, amplifying, and oscillatory families.
\end{itemize}

Together, RPM and Fractals provide the complete dynamic semantics needed to evaluate TKS expressions, diagnose failures, and analyze transformation stability.

\section{Notation and Conventions}
\label{sec:notation}

Throughout this manual, we employ consistent notation:

\begin{notation}[Symbol Classes]
\label{not:symbols}
The following symbol classes are used:
\begin{center}
\begin{tabular}{lll}
\toprule
\textbf{Class} & \textbf{Notation} & \textbf{Domain} \\
\midrule
Worlds & $A, B, C, D$ & $\Worlds = \{A, B, C, D\}$ \\
Noetics & $\noe{0}, \noe{1}, \ldots, \noe{9}$ & $\Noetics = \{\noe{k} : k \in \{0,\ldots,9\}\}$ \\
Foundations & $F_1, F_2, \ldots, F_7$ & $\Foundations = \{F_m : m \in \{1,\ldots,7\}\}$ \\
Elements & $Xn$ where $X \in \Worlds$, $n \in \{1,\ldots,10\}$ & $\Elements = \Worlds \times \{1,\ldots,10\}$ \\
Ideas & $\Ideas$ & Space of all Ideas \\
Fractals & $\Frac, \Psi, \Omega$ & $\FracSet$ \\
Acquisitions & $A0, D_m, W_m, P_m$ & $\Acquisitions$ \\
\bottomrule
\end{tabular}
\end{center}
\end{notation}

\begin{notation}[Operations]
\label{not:operations}
\begin{center}
\begin{tabular}{lll}
\toprule
\textbf{Operation} & \textbf{Symbol} & \textbf{Meaning} \\
\midrule
Tootra-Addition & $\Tadd$ & Idea union/co-presence \\
Tootra-Subtraction & $\Tsub$ & Idea removal \\
Tootra-Multiplication & $\Tmul$ & Idea interaction/scaling \\
Tootra-Division & $\Tdiv$ & Idea decomposition \\
World Association & $\oplus$ & Link expression to world \\
Composition & $\circ$ & Sequential application \\
Dependency & $\to$ & Prerequisite relation \\
Fractal & $\langle X : Y : Z \rangle$ & Noetic fractal chain \\
Semantic bracket & $\sem{\cdot}$ & Denotational meaning \\
RPM return & $\rpmret$ & Monad unit \\
RPM bind & $\rpmbind$ & Kleisli composition \\
Big-step eval & $\bigstep$ & Evaluates to \\
Small-step & $\smallstep$ & Single reduction step \\
\bottomrule
\end{tabular}
\end{center}
\end{notation}

\chapter{Canonical Ontology}
\label{ch:ontology}

This chapter establishes the foundational ontology of TKS: the primordial categories, the Four Worlds, the 40 Elements, the 10 Noetic operators, and the 7 Foundations. All definitions in this chapter are canonical from v6.0 and preserved unchanged.

\section{The Primordial Duality: Mind and Idea}

\begin{axiom}[Ontological Foundation]
\label{ax:ontological-foundation}
The TKS system posits two primordial categories from which all else derives:
\begin{align}
\textbf{MIND} &: \text{The operator that processes, stores, and transforms} \\
\textbf{IDEA} &: \text{The operand---all that can be represented or structured}
\end{align}
\end{axiom}

\begin{definition}[Idea Space]
\label{def:idea-space}
The \textbf{Idea space} is the collection of all Ideas:
\[
\Ideas = \{\text{all possible Ideas}\}
\]
An Idea encompasses spiritual, mental, emotional, or physical content, including any composite expression built from Elements, Foundations, and Worlds.
\end{definition}

\begin{definition}[Mind as Operator Class]
\label{def:mind-class}
Mind is an operator class $M$ with instantiations across four worlds:
\[
M = \{M_A, M_B, M_C, M_D\}
\]
where:
\begin{itemize}
  \item $M_A \cong A1$: Divine Mind (First Cause Awareness)
  \item $M_B \cong B1$: Mental Mind (Ego, Memory, Analytical Intelligence)
  \item $M_C \cong C1$: Emotional Mind (Emotional Awareness, EQ)
  \item $M_D \cong D1$: Physical Mind (Brain, Hardware, Biological System)
\end{itemize}
\end{definition}

\section{The Four Worlds}

\begin{definition}[World Set]
\label{def:world-set}
The \textbf{World set} is:
\[
\Worlds = \{A, B, C, D\}
\]
with total ordering by metaphysical subtlety:
\[
A \succ B \succ C \succ D
\]
\end{definition}

\begin{definition}[World Ordinal Function]
\label{def:world-ord}
The \textbf{ordinal function} assigns a natural number to each world:
\[
\ord : \Worlds \to \mathbb{N}
\]
\[
\ord(A) = 0, \quad \ord(B) = 1, \quad \ord(C) = 2, \quad \ord(D) = 3
\]
\end{definition}

\begin{definition}[World Distance]
\label{def:world-distance}
The \textbf{distance} between worlds $W_1$ and $W_2$ is:
\[
d(W_1, W_2) = |\ord(W_1) - \ord(W_2)|
\]
\end{definition}

\begin{table}[h]
\centering
\caption{World Signatures}
\label{tab:world-signatures}
\begin{tabular}{llllll}
\toprule
\textbf{World} & \textbf{Domain} & \textbf{Substrate} & \textbf{Kabbalistic} & \textbf{Mind} & \textbf{Idea} \\
\midrule
$A$ & Spiritual & Aether & Atziluth & $A1$ & $A10$ \\
$B$ & Mental & Thought & Briah & $B1$ & $B10$ \\
$C$ & Emotional & Feeling & Yetzirah & $C1$ & $C10$ \\
$D$ & Physical & Matter & Assiah & $D1$ & $D10$ \\
\bottomrule
\end{tabular}
\end{table}

\section{The 40 Elements}

\begin{definition}[Element Space]
\label{def:element-space}
The \textbf{Element space} is the Cartesian product:
\[
\Elements = \Worlds \times \{1, 2, 3, 4, 5, 6, 7, 8, 9, 10\}
\]
yielding $|\Elements| = 4 \times 10 = 40$ elements.
\end{definition}

\begin{theorem}[Element Tensor Structure]
\label{thm:element-tensor}
The 40 Elements form a $4 \times 10$ tensor grid:

\[
\begin{array}{c|cccccccccc}
 & 1 & 2 & 3 & 4 & 5 & 6 & 7 & 8 & 9 & 10 \\
\hline
A & A1 & A2 & A3 & A4 & A5 & A6 & A7 & A8 & A9 & A10 \\
B & B1 & B2 & B3 & B4 & B5 & B6 & B7 & B8 & B9 & B10 \\
C & C1 & C2 & C3 & C4 & C5 & C6 & C7 & C8 & C9 & C10 \\
D & D1 & D2 & D3 & D4 & D5 & D6 & D7 & D8 & D9 & D10 \\
\end{array}
\]
\end{theorem}

\subsection{Element Semantics by Mode}

Each mode $n \in \{1, \ldots, 10\}$ carries a consistent meaning across all worlds:

\begin{center}
\begin{tabular}{cl}
\toprule
\textbf{Mode} & \textbf{Semantic Role} \\
\midrule
1 & Mind / Awareness / Processing \\
2 & Positive / Attraction / Union \\
3 & Negative / Repulsion / Separation \\
4 & Vibration / Frequency / Energy \\
5 & Female / Receptive / Nurturing \\
6 & Male / Projective / Structuring \\
7 & Rhythm / Cycles / Periodicity \\
8 & Above / Inner / Esoteric / Causal \\
9 & Below / Outer / Exoteric / Effect \\
10 & Idea / Pure Form / Potential \\
\bottomrule
\end{tabular}
\end{center}

\section{The Ten Noetic Operators}

\begin{definition}[Noetic Operator Space]
\label{def:noetic-space}
The \textbf{Noetic operator space} is:
\[
\Noetics = \{\noe{0}, \noe{1}, \noe{2}, \noe{3}, \noe{4}, \noe{5}, \noe{6}, \noe{7}, \noe{8}, \noe{9}\}
\]
Each Noetic is an endofunction on the Idea space:
\[
\noe{k} : \Ideas \to \Ideas
\]
\end{definition}

\begin{table}[h]
\centering
\caption{Noetic Operators}
\label{tab:noetics}
\begin{tabular}{clll}
\toprule
\textbf{Index} & \textbf{Symbol} & \textbf{Name} & \textbf{Function} \\
\midrule
0 & $\noe{0}$ & Idea & Neutral form, identity \\
1 & $\noe{1}$ & Mind & Attention, awareness, processing \\
2 & $\noe{2}$ & Positive & Attraction, affirmation, union \\
3 & $\noe{3}$ & Negative & Repulsion, negation, separation \\
4 & $\noe{4}$ & Vibration & Amplitude/frequency modulation \\
5 & $\noe{5}$ & Female & Receptive structuring, internalization \\
6 & $\noe{6}$ & Male & Projective structuring, externalization \\
7 & $\noe{7}$ & Rhythm & Periodicity, repetition, cycles \\
8 & $\noe{8}$ & Above & Inner, higher, esoteric, causal \\
9 & $\noe{9}$ & Below & Outer, lower, exoteric, effect \\
\bottomrule
\end{tabular}
\end{table}

\subsection{The Mirror Principle}

\begin{definition}[Noetic Mirror Pairs]
\label{def:mirror-pairs}
Noetics form \textbf{mirror pairs} that sum to 9:
\begin{center}
\begin{tabular}{ccc}
\toprule
\textbf{Pair} & \textbf{Sum} & \textbf{Interpretation} \\
\midrule
$\noe{1} \leftrightarrow \noe{8}$ & $1 + 8 = 9$ & Mind $\leftrightarrow$ Above \\
$\noe{2} \leftrightarrow \noe{7}$ & $2 + 7 = 9$ & Positive $\leftrightarrow$ Rhythm \\
$\noe{3} \leftrightarrow \noe{6}$ & $3 + 6 = 9$ & Negative $\leftrightarrow$ Male \\
$\noe{4} \leftrightarrow \noe{5}$ & $4 + 5 = 9$ & Vibration $\leftrightarrow$ Female \\
\bottomrule
\end{tabular}
\end{center}
\end{definition}

\section{The Seven Foundations}

\begin{definition}[Foundation Space]
\label{def:foundation-space}
The \textbf{Foundation space} is:
\[
\Foundations = \{F_1, F_2, F_3, F_4, F_5, F_6, F_7\}
\]
\end{definition}

\begin{table}[h]
\centering
\caption{The Seven Foundations}
\label{tab:foundations}
\begin{tabular}{clll}
\toprule
$m$ & \textbf{Foundation} & \textbf{Core Meaning} & \textbf{Life Domain} \\
\midrule
1 & Unity & Coherence, integration & Spiritual wholeness \\
2 & Wisdom & Knowledge, understanding & Truth and learning \\
3 & Life & Vitality, health & Biological survival \\
4 & Companionship & Connection, love & Relationships \\
5 & Power & Influence, agency & Authority and action \\
6 & Material & Resources, possessions & Economic domain \\
7 & Lust & Sex, reproduction & Procreation \\
\bottomrule
\end{tabular}
\end{table}

\begin{definition}[Sub-Foundation Structure]
\label{def:subfoundation}
Each Foundation has four world-manifestations:
\[
F_m = \{F_{ma}, F_{mb}, F_{mc}, F_{md}\}
\]
where the subscript indicates the world ($a$ = Spiritual, $b$ = Mental, $c$ = Emotional, $d$ = Physical).
\end{definition}

\begin{definition}[Sub-Foundation Space]
\label{def:subfoundation-space}
The complete \textbf{Sub-Foundation space} is:
\[
\Foundations^* = \Foundations \times \{a, b, c, d\}
\]
with $|\Foundations^*| = 7 \times 4 = 28$ sub-foundations.
\end{definition}

\begin{table}[ht]
\centering
\caption{The 28 Sub-Foundations}
\label{tab:28-subfoundations}
\begin{tabular}{c|cccc}
\toprule
\textbf{Foundation} & \textbf{a (Spiritual)} & \textbf{b (Mental)} & \textbf{c (Emotional)} & \textbf{d (Physical)} \\
\midrule
$F_1$ (Unity) & $1_a$ & $1_b$ & $1_c$ & $1_d$ \\
$F_2$ (Wisdom) & $2_a$ & $2_b$ & $2_c$ & $2_d$ \\
$F_3$ (Life) & $3_a$ & $3_b$ & $3_c$ & $3_d$ \\
$F_4$ (Companionship) & $4_a$ & $4_b$ & $4_c$ & $4_d$ \\
$F_5$ (Power) & $5_a$ & $5_b$ & $5_c$ & $5_d$ \\
$F_6$ (Material) & $6_a$ & $6_b$ & $6_c$ & $6_d$ \\
$F_7$ (Lust) & $7_a$ & $7_b$ & $7_c$ & $7_d$ \\
\bottomrule
\end{tabular}
\end{table}

\begin{proposition}[Sub-Foundation Lattice]
\label{prop:subfoundation-lattice}
The sub-foundation space $\Foundations^*$ forms a bounded distributive lattice $(\Foundations^*, \sqcup, \sqcap, \bot, \top)$ where:
\begin{enumerate}[label=(\roman*)]
  \item Join $\sqcup$ combines foundation development
  \item Meet $\sqcap$ extracts common foundation
  \item $\bot$ = no foundation active
  \item $\top$ = all 28 sub-foundations fully developed
\end{enumerate}
\end{proposition}

\begin{example}[Sub-Foundation Interpretations]
\label{ex:subfoundation-interpretations}
\begin{itemize}
  \item $3_d$ (Physical Life): biological energy, health, stamina
  \item $5_b$ (Mental Power): influence through ideas, persuasion
  \item $7_c$ (Emotional Lust): desire, attraction, bonding
  \item $2_a$ (Spiritual Wisdom): divine knowledge, archetypal truth
\end{itemize}
\end{example}

\section{The Acquisition Space}

\begin{definition}[Acquisition Space]
\label{def:acq-space}
The \textbf{Acquisition space} is:
\[
\Acquisitions = \{A0\} \cup \{D_m : m \in \{1,\ldots,7\}\} \cup \{W_m : m \in \{1,\ldots,7\}\} \cup \{P_m : m \in \{1,\ldots,7\}\}
\]
with $|\Acquisitions| = 1 + 7 + 7 + 7 = 22$.

\begin{itemize}
  \item $A0$: Pure Desire (root intentional vector)
  \item $D_m$: Desire for Foundation $m$
  \item $W_m$: Wisdom about Foundation $m$
  \item $P_m$: Power to achieve Foundation $m$
\end{itemize}
\end{definition}

\begin{definition}[Acquisition State]
\label{def:acq-state}
An \textbf{Acquisition State} is a function:
\[
\AcqState : \Acquisitions \to \{\mathtt{true}, \mathtt{false}\}
\]
indicating which acquisitions are currently satisfied.
\end{definition}

%% ============================================================================
%% PART II: ALGEBRAIC STRUCTURES
%% [CONTENT PRESERVED FROM v6.1 - TO BE COPIED]
%% ============================================================================
\part{Algebraic Structures}

\chapter{Noetic Operator Algebra}
\label{ch:noetic-algebra}

This chapter provides a complete specification of the algebraic structure of Noetic operators, including the composition table, structural properties, and eigenmodes. All definitions are preserved from v6.0.

\section{Composition as Binary Operation}

\begin{definition}[Noetic Composition]
\label{def:noetic-composition}
The \textbf{composition operation} on Noetics is:
\[
\circ : \Noetics \times \Noetics \to \mathrm{End}(\Ideas)
\]
defined by:
\[
(\noe{j} \circ \noe{i})(X) = \noe{j}(\noe{i}(X))
\]
for all $X \in \Ideas$.
\end{definition}

\begin{remark}
The notation $\noe{j} \circ \noe{i}$ means ``apply $\noe{i}$ first, then $\noe{j}$.'' This follows the standard convention for function composition.
\end{remark}

\section{Fundamental Composition Axioms}

\begin{axiom}[Associativity of Composition]
\label{ax:assoc-comp}
For all $\noe{i}, \noe{j}, \noe{k} \in \Noetics$:
\[
\noe{k} \circ (\noe{j} \circ \noe{i}) = (\noe{k} \circ \noe{j}) \circ \noe{i}
\]
\end{axiom}

\begin{axiom}[Identity Element]
\label{ax:identity}
The Noetic $\noe{0}$ (Idea) serves as the identity element:
\[
\noe{0} \circ \noe{k} = \noe{k} \circ \noe{0} = \noe{k} \quad \forall k \in \{0,\ldots,9\}
\]
\end{axiom}

\begin{theorem}[Monoid Structure]
\label{thm:monoid}
The structure $(\Noetics, \circ, \noe{0})$ forms a monoid.
\end{theorem}

\begin{proof}
We verify the monoid axioms:
\begin{enumerate}
  \item \textbf{Closure}: For any $\noe{i}, \noe{j} \in \Noetics$, the composition $\noe{j} \circ \noe{i}$ is an endofunction on $\Ideas$. By the closure axiom of TKS, compositions of Noetics yield valid operators in the extended closure $\overline{\Noetics}$.
  \item \textbf{Associativity}: By Axiom~\ref{ax:assoc-comp}.
  \item \textbf{Identity}: By Axiom~\ref{ax:identity}, $\noe{0}$ is the identity.
\end{enumerate}
\end{proof}

\section{The Complete 10$\times$10 Noetic Composition Table}

\begin{definition}[Composition Notation]
\label{def:comp-notation}
We use the following notational conventions:
\begin{itemize}
  \item $\noe{ij}$ denotes the compound operator $\noe{j} \circ \noe{i}$
  \item $\noe{i}^2$ denotes the self-composition $\noe{i} \circ \noe{i}$
  \item $\approx \noe{0}$ indicates approximate neutralization (returns toward identity)
\end{itemize}
\end{definition}

\begin{table}[h]
\centering
\caption{Complete 10$\times$10 Noetic Composition Table}
\label{tab:composition}
\small
\begin{tabular}{c|cccccccccc}
$\circ$ & $\noe{0}$ & $\noe{1}$ & $\noe{2}$ & $\noe{3}$ & $\noe{4}$ & $\noe{5}$ & $\noe{6}$ & $\noe{7}$ & $\noe{8}$ & $\noe{9}$ \\
\hline
$\noe{0}$ & $\noe{0}$ & $\noe{1}$ & $\noe{2}$ & $\noe{3}$ & $\noe{4}$ & $\noe{5}$ & $\noe{6}$ & $\noe{7}$ & $\noe{8}$ & $\noe{9}$ \\
$\noe{1}$ & $\noe{1}$ & $\noe{1}^2$ & $\noe{12}$ & $\noe{13}$ & $\noe{14}$ & $\noe{15}$ & $\noe{16}$ & $\noe{17}$ & $\noe{18}$ & $\noe{19}$ \\
$\noe{2}$ & $\noe{2}$ & $\noe{21}$ & $\noe{2}^2$ & $\approx\noe{0}$ & $\noe{24}$ & $\noe{25}$ & $\noe{26}$ & $\noe{27}$ & $\noe{28}$ & $\noe{29}$ \\
$\noe{3}$ & $\noe{3}$ & $\noe{31}$ & $\approx\noe{0}$ & $\noe{3}^2$ & $\noe{34}$ & $\noe{35}$ & $\noe{36}$ & $\noe{37}$ & $\noe{38}$ & $\noe{39}$ \\
$\noe{4}$ & $\noe{4}$ & $\noe{41}$ & $\noe{42}$ & $\noe{43}$ & $\noe{4}^2$ & $\noe{45}$ & $\noe{46}$ & $\noe{47}$ & $\noe{48}$ & $\noe{49}$ \\
$\noe{5}$ & $\noe{5}$ & $\noe{51}$ & $\noe{52}$ & $\noe{53}$ & $\noe{54}$ & $\noe{5}^2$ & $\approx\noe{0}$ & $\noe{57}$ & $\noe{58}$ & $\noe{59}$ \\
$\noe{6}$ & $\noe{6}$ & $\noe{61}$ & $\noe{62}$ & $\noe{63}$ & $\noe{64}$ & $\approx\noe{0}$ & $\noe{6}^2$ & $\noe{67}$ & $\noe{68}$ & $\noe{69}$ \\
$\noe{7}$ & $\noe{7}$ & $\noe{71}$ & $\noe{72}$ & $\noe{73}$ & $\noe{74}$ & $\noe{75}$ & $\noe{76}$ & $\noe{7}^2$ & $\noe{78}$ & $\noe{79}$ \\
$\noe{8}$ & $\noe{8}$ & $\noe{81}$ & $\noe{82}$ & $\noe{83}$ & $\noe{84}$ & $\noe{85}$ & $\noe{86}$ & $\noe{87}$ & $\noe{8}^2$ & $\approx\noe{0}$ \\
$\noe{9}$ & $\noe{9}$ & $\noe{91}$ & $\noe{92}$ & $\noe{93}$ & $\noe{94}$ & $\noe{95}$ & $\noe{96}$ & $\noe{97}$ & $\approx\noe{0}$ & $\noe{9}^2$ \\
\end{tabular}
\end{table}

\section{Dualities and Inversions}

\begin{definition}[Dual Noetics]
\label{def:dual-noetics}
The three \textbf{fundamental dualities} are:
\begin{align}
\text{Duality}_1 &: \noe{2} \leftrightarrow \noe{3} && \text{(Positive } \leftrightarrow \text{ Negative)} \\
\text{Duality}_2 &: \noe{5} \leftrightarrow \noe{6} && \text{(Female } \leftrightarrow \text{ Male)} \\
\text{Duality}_3 &: \noe{8} \leftrightarrow \noe{9} && \text{(Above } \leftrightarrow \text{ Below)}
\end{align}
\end{definition}

\begin{axiom}[Pseudo-Inverse Relations]
\label{ax:pseudo-inverse}
For dual pairs, we define pseudo-inverses:
\begin{align}
\noe{2}^{-1} &:= \noe{3} & \noe{3}^{-1} &:= \noe{2} \\
\noe{5}^{-1} &:= \noe{6} & \noe{6}^{-1} &:= \noe{5} \\
\noe{8}^{-1} &:= \noe{9} & \noe{9}^{-1} &:= \noe{8}
\end{align}
\end{axiom}

\begin{theorem}[Duality Neutralization]
\label{thm:duality-neutral}
For dual pairs $(\noe{a}, \noe{b})$:
\[
\noe{b} \circ \noe{a} \approx \noe{0}
\]
The composition returns toward the neutral/potential state.
\end{theorem}

\begin{proof}
We verify for each dual pair:
\begin{enumerate}
  \item $\noe{3} \circ \noe{2}$: Attraction followed by repulsion nullifies the directional component, leaving undifferentiated potential.
  \item $\noe{6} \circ \noe{5}$: Reception followed by projection restructures content back toward potential.
  \item $\noe{9} \circ \noe{8}$: Elevation followed by grounding returns content toward neutral manifestation level.
\end{enumerate}
The approximate equality $\approx$ indicates that while the result may not be exactly $\noe{0}$ in all contexts, it is semantically equivalent to returning toward the identity operation.
\end{proof}

\section{Commutators and Anti-Commutators}

\begin{definition}[Noetic Commutator]
\label{def:commutator}
For Noetics $\noe{i}, \noe{j}$, the \textbf{commutator} is:
\[
[\noe{i}, \noe{j}] := \noe{i} \circ \noe{j} - \noe{j} \circ \noe{i}
\]
where subtraction is interpreted as the operator difference.
\end{definition}

\begin{theorem}[Key Commutator Relations]
\label{thm:commutators}
The following commutator relations hold:
\begin{enumerate}
  \item $[\noe{0}, \noe{k}] = 0$ for all $k$ (identity commutes with everything)
  \item $[\noe{2}, \noe{3}] \neq 0$ (Positive-Negative order matters)
  \item $[\noe{5}, \noe{6}] \neq 0$ (Female-Male order matters)
  \item $[\noe{8}, \noe{9}] \neq 0$ (Above-Below order matters)
  \item $[\noe{4}, \noe{7}] \approx 0$ (Vibration and Rhythm approximately commute)
\end{enumerate}
\end{theorem}

\section{Eigenmodes and Stability}

\begin{definition}[Noetic Eigenstate]
\label{def:eigenstate}
An Idea $X \in \Ideas$ is an \textbf{eigenstate} of Noetic $\noe{k}$ with eigenvalue $\lambda$ if:
\[
\noe{k}(X) = \lambda X
\]
\end{definition}

\begin{theorem}[Element-Noetic Correspondence]
\label{thm:element-noetic}
Each Element $Xn$ (where $X \in \Worlds$ and $n \in \{1,\ldots,10\}$) is a stable eigenstate of its corresponding Noetic $\noe{n}$:
\[
\noe{n}(Xn) = Xn \quad (\lambda = 1)
\]
\end{theorem}

\begin{proof}
Each Element of mode $n$ is intrinsically aligned with Noetic $\noe{n}$. Applying $\noe{n}$ to an already-$n$-aligned Element preserves its form.
\end{proof}

\chapter{Algebraic Structures of Tootra Arithmetic}
\label{ch:tootra-arithmetic}

This chapter provides a complete formalization of Tootra arithmetic on the Idea space.

\section{The Idea Space as a Mathematical Structure}

\begin{definition}[Idea Space Structure]
\label{def:idea-structure}
The \textbf{Idea space} $\Ideas$ is defined as a \textbf{graded multiset algebra}:
\[
\Ideas = \bigcup_{W \in \Worlds} \mathcal{M}(\mathcal{P}_W)
\]
where $\mathcal{M}(\mathcal{P}_W)$ denotes the space of finite multisets over a property space $\mathcal{P}_W$ for world $W$.
\end{definition}

\begin{definition}[Property Space]
\label{def:property-space}
For each world $W$, the \textbf{property space} $\mathcal{P}_W$ is a set of atomic properties:
\[
\mathcal{P}_W = \{p_{W,i} : i \in I_W\}
\]
An Idea in world $W$ is a multiset $X : \mathcal{P}_W \to \mathbb{N}$ assigning a non-negative multiplicity to each property.
\end{definition}

\section{Tootra-Addition ($\Tadd$)}

\begin{definition}[Tootra-Addition]
\label{def:tootra-add}
For Ideas $X, Y \in \Ideas$, the \textbf{Tootra-Addition} is:
\[
X \Tadd Y = X \uplus Y
\]
where $\uplus$ denotes multiset union (summing multiplicities).

Formally, for each property $p$:
\[
(X \Tadd Y)(p) = X(p) + Y(p)
\]
\end{definition}

\begin{theorem}[Properties of Tootra-Addition]
\label{thm:tadd-properties}
Tootra-Addition satisfies:
\begin{enumerate}
  \item \textbf{Commutativity}: $X \Tadd Y = Y \Tadd X$
  \item \textbf{Associativity}: $(X \Tadd Y) \Tadd Z = X \Tadd (Y \Tadd Z)$
  \item \textbf{Identity}: $X \Tadd \varnothing = X$ (where $\varnothing$ is the empty multiset)
\end{enumerate}
\end{theorem}

\begin{proof}
\begin{enumerate}
  \item \textbf{Commutativity}: $(X \Tadd Y)(p) = X(p) + Y(p) = Y(p) + X(p) = (Y \Tadd X)(p)$.
  \item \textbf{Associativity}: $((X \Tadd Y) \Tadd Z)(p) = (X(p) + Y(p)) + Z(p) = X(p) + (Y(p) + Z(p))$.
  \item \textbf{Identity}: $(X \Tadd \varnothing)(p) = X(p) + 0 = X(p)$.
\end{enumerate}
\end{proof}

\begin{corollary}[Abelian Monoid Structure]
\label{cor:tadd-monoid}
The structure $(\Ideas, \Tadd, \varnothing)$ forms an \textbf{abelian monoid}.
\end{corollary}

\section{Tootra-Subtraction ($\Tsub$)}

\begin{definition}[Tootra-Subtraction]
\label{def:tootra-sub}
For Ideas $X, Y \in \Ideas$, the \textbf{Tootra-Subtraction} is:
\[
X \Tsub Y = X \setminus_m Y
\]
where $\setminus_m$ denotes multiset difference (subtracting multiplicities, floored at 0).

Formally:
\[
(X \Tsub Y)(p) = \max(0, X(p) - Y(p))
\]
\end{definition}

\begin{theorem}[Properties of Tootra-Subtraction]
\label{thm:tsub-properties}
Tootra-Subtraction satisfies:
\begin{enumerate}
  \item \textbf{Non-Commutativity}: $X \Tsub Y \neq Y \Tsub X$ in general
  \item \textbf{Right Identity}: $X \Tsub \varnothing = X$
  \item \textbf{Nullification}: $X \Tsub X = \varnothing$
\end{enumerate}
\end{theorem}

\section{Tootra-Multiplication ($\Tmul$)}

\begin{definition}[Tootra-Multiplication]
\label{def:tootra-mul}
Under the simplified (Hadamard) form:
\[
(X \Tmul Y)(p) = X(p) \cdot Y(p)
\]
This is the component-wise product.
\end{definition}

\begin{theorem}[Properties of Tootra-Multiplication]
\label{thm:tmul-properties}
Under simplified Tootra-Multiplication:
\begin{enumerate}
  \item \textbf{Commutativity}: $X \Tmul Y = Y \Tmul X$
  \item \textbf{Associativity}: $(X \Tmul Y) \Tmul Z = X \Tmul (Y \Tmul Z)$
  \item \textbf{Identity}: $X \Tmul \mathbf{1} = X$ where $\mathbf{1}(p) = 1$ for all $p$
  \item \textbf{Annihilation}: $X \Tmul \varnothing = \varnothing$
\end{enumerate}
\end{theorem}

\section{Tootra-Division ($\Tdiv$)}

\begin{definition}[Tootra-Division]
\label{def:tootra-div}
For Ideas $X, Y \in \Ideas$ where $Y \neq \varnothing$:
\[
(X \Tdiv Y)(p) = \begin{cases}
\lfloor X(p) / Y(p) \rfloor & \text{if } Y(p) > 0 \\
X(p) & \text{if } Y(p) = 0
\end{cases}
\]
\end{definition}

\section{Algebraic Structure Summary}

\begin{theorem}[Tootra Semiring Structure]
\label{thm:semiring}
Under simplified operations, $(\Ideas, \Tadd, \Tmul, \varnothing, \mathbf{1})$ forms a \textbf{commutative semiring}.
\end{theorem}

\begin{proof}
We verify the semiring axioms:
\begin{enumerate}
  \item $(\Ideas, \Tadd, \varnothing)$ is a commutative monoid.
  \item $(\Ideas, \Tmul, \mathbf{1})$ is a commutative monoid.
  \item \textbf{Distributivity}: $(X \Tadd Y) \Tmul Z = (X \Tmul Z) \Tadd (Y \Tmul Z)$
  \item \textbf{Annihilation}: $X \Tmul \varnothing = \varnothing$
\end{enumerate}
\end{proof}

%% ============================================================================
%% PART III: CATEGORICAL FRAMEWORK
%% [CONTENT PRESERVED FROM v6.1 - TO BE COPIED]
%% ============================================================================
\part{Categorical Framework}

\chapter{The Category Noetica}
\label{ch:category-noetica}

This chapter provides a rigorous category-theoretic foundation for TKS.

\section{Definition of Category Noetica}

\begin{definition}[Category Noetica]
\label{def:category-noetica}
The category $\Noetica$ is defined by:

\textbf{Objects:}
\[
\Ob(\Noetica) = \Ideas \cup \Worlds \cup \Elements \cup \Foundations \cup \Acquisitions
\]

\textbf{Morphisms:}
\[
\Hom(\Noetica) = \Noetics \cup \{\oplus_W, \ominus_W : W \in \Worlds\}
\]

For objects $X, Y \in \Ob(\Noetica)$:
\[
\Hom(X, Y) = \{f : X \to Y \mid f \text{ is a valid TKS transformation}\}
\]

\textbf{Identity:}
\[
\id_X = \noe{0} : X \to X
\]

\textbf{Composition:}
For morphisms $f : X \to Y$ and $g : Y \to Z$:
\[
(g \circ f) : X \to Z, \quad (g \circ f)(x) = g(f(x))
\]
\end{definition}

\section{Verification of Category Axioms}

\begin{theorem}[Noetica is a Category]
\label{thm:noetica-category}
The structure $\Noetica$ satisfies all category axioms.
\end{theorem}

\begin{proof}
We verify each axiom:

\textbf{(C1) Identity Law:}
For any morphism $f : X \to Y$:
\begin{align}
f \circ \id_X &= f \circ \noe{0} = f \\
\id_Y \circ f &= \noe{0} \circ f = f
\end{align}
Both equalities hold because $\noe{0}(x) = x$ for all $x$ (Axiom~\ref{ax:identity}).

\textbf{(C2) Associativity:}
For morphisms $f : X \to Y$, $g : Y \to Z$, $h : Z \to W$:
\begin{align}
(h \circ (g \circ f))(x) &= h((g \circ f)(x)) = h(g(f(x))) \\
((h \circ g) \circ f)(x) &= (h \circ g)(f(x)) = h(g(f(x)))
\end{align}
Thus $h \circ (g \circ f) = (h \circ g) \circ f$.

\textbf{(C3) Composition Closure:}
For any $f \in \Hom(X, Y)$ and $g \in \Hom(Y, Z)$, the composite $g \circ f$ is in $\Hom(X, Z)$ by closure of Noetic composition.
\end{proof}

\section{World Subcategories}

\begin{definition}[World Subcategory]
\label{def:world-subcat}
For each World $W \in \{A, B, C, D\}$, define the subcategory $\Noetica_W$:
\begin{align}
\Ob(\Noetica_W) &= \{X \in \Ob(\Noetica) : \mathrm{World}(X) = W\} \\
\Hom(\Noetica_W) &= \{f \in \Hom(\Noetica) : f \text{ preserves World } W\}
\end{align}
\end{definition}

\begin{proposition}[World Subcategories are Categories]
\label{prop:world-subcat}
Each $\Noetica_W$ is a valid subcategory of $\Noetica$.
\end{proposition}

\chapter{The ACBE Functor}
\label{ch:acbe-functor}

The ACBE (Above-Cause-Below-Effect) transformation is the primary manifestation functor in TKS.

\section{Functor Definition}

\begin{definition}[ACBE Functor]
\label{def:acbe-functor}
The \textbf{ACBE functor} is:
\[
\mathrm{ACBE} : \Noetica_A \to \Noetica_D
\]

\textbf{Object Mapping:}
For each $An \in \Ob(\Noetica_A)$:
\[
\mathrm{ACBE}(An) = Dn
\]

\textbf{Morphism Mapping:}
For morphism $f : X \to Y$ in $\Noetica_A$:
\[
\mathrm{ACBE}(f) : \mathrm{ACBE}(X) \to \mathrm{ACBE}(Y)
\]
\[
\mathrm{ACBE}(\noe{k}) = \noe{k}
\]
Noetics are preserved across worlds.
\end{definition}

\section{Proof of Functoriality}

\begin{theorem}[ACBE is a Functor]
\label{thm:acbe-functor}
ACBE preserves identity and composition, hence is a valid functor.
\end{theorem}

\begin{proof}
\textbf{(F1) Identity Preservation:}
\[
\mathrm{ACBE}(\id_X) = \mathrm{ACBE}(\noe{0}) = \noe{0} = \id_{\mathrm{ACBE}(X)}
\]

\textbf{(F2) Composition Preservation:}
For $f : X \to Y$ and $g : Y \to Z$ in $\Noetica_A$:
\[
\mathrm{ACBE}(g \circ f) = \mathrm{ACBE}(g) \circ \mathrm{ACBE}(f)
\]
Since $\mathrm{ACBE}(\noe{j} \circ \noe{i}) = \noe{j} \circ \noe{i} = \mathrm{ACBE}(\noe{j}) \circ \mathrm{ACBE}(\noe{i})$.
\end{proof}

\section{The Cascade Decomposition}

\begin{theorem}[ACBE Factorization]
\label{thm:acbe-factor}
ACBE factors through intermediate categories:
\[
\mathrm{ACBE} = F_D \circ F_C \circ F_B
\]
where:
\begin{align}
F_B &: \Noetica_A \to \Noetica_B && \text{(Spiritual } \to \text{ Mental)} \\
F_C &: \Noetica_B \to \Noetica_C && \text{(Mental } \to \text{ Emotional)} \\
F_D &: \Noetica_C \to \Noetica_D && \text{(Emotional } \to \text{ Physical)}
\end{align}
\end{theorem}

\begin{proof}
Each $F_W$ is defined by $F_B(An) = Bn$, $F_C(Bn) = Cn$, $F_D(Cn) = Dn$.
Then $(F_D \circ F_C \circ F_B)(An) = Dn = \mathrm{ACBE}(An)$.
\end{proof}

%% ============================================================================
%% PART IV: RPM AND FRACTAL CALCULUS (EXPANDED IN v7.0)
%% ============================================================================
\part{RPM Monad Theory and Noetic Fractal Calculus}

\chapter{The RPM Monad: Complete Theory}
\label{ch:rpm-monad}

\begin{tcolorbox}[colback=green!5!white,colframe=green!75!black,title=\vexp]
This chapter expands the v6.1 RPM monad formalization with complete category-theoretic treatment, Kleisli category structure, effect system interpretation, and multi-goal parallel structures.
\end{tcolorbox}

\section{Motivation: Prerequisites as Computational Effects}

In TKS, not every transformation can be performed unconditionally. The RPM captures the constraint that certain acquisitions (Desire, Wisdom, Power) must be satisfied before outcomes can manifest. This is naturally modeled as a \emph{computational effect}: evaluation proceeds through a chain of prerequisite checks, and failure at any point aborts the computation.

\begin{definition}[Prerequisite Operators]
\label{def:prereq-ops}
The \textbf{Prerequisite operator set} is:
\[
\Pi = \{D, W, P\}
\]
where:
\begin{itemize}
  \item $D$: Desire --- initiates and directs
  \item $W$: Wisdom --- models and validates
  \item $P$: Power --- executes and manifests
\end{itemize}
\end{definition}

\begin{theorem}[Strict Ordering]
\label{thm:prereq-order}
The prerequisite operators have a strict total ordering:
\[
D \prec W \prec P
\]
Desire must precede Wisdom must precede Power.
\end{theorem}

\section{The RPM Type Constructor}

\begin{definition}[RPM Type Constructor]
\label{def:rpm-type}
The \textbf{RPM type constructor} is:
\[
\RPM : \mathcal{D} \to \mathcal{D}_\RPM
\]
where $\mathcal{D}$ is the semantic domain and $\mathcal{D}_\RPM$ is the domain of RPM-wrapped computations.

Concretely:
\[
\RPM(A) = \AcqState \to (A + \Failed) \times \AcqState
\]
An RPM computation takes an acquisition state $\AcqState$, and returns either a value of type $A$ together with an updated state, or a failure indicator with the failure origin.
\end{definition}

\begin{definition}[RPM Result Type]
\label{def:rpm-result}
The result of an RPM computation is:
\[
\mathsf{Result}_A = \mathsf{Success}(a : A, \sigma' : \AcqState) \mid \mathsf{Failure}(o : \Acquisitions)
\]
where $o$ is the \textbf{failure origin}---the acquisition that was not satisfied.
\end{definition}

\section{Monad Structure}

\begin{definition}[RPM Unit (Return)]
\label{def:rpm-unit}
The \textbf{unit} (return) operation $\eta : A \to \RPM(A)$ is:
\[
\rpmret(a) = \lambda \sigma. \, (\mathsf{Success}(a), \sigma)
\]
It injects a pure value into the RPM monad without modifying state.
\end{definition}

\begin{definition}[RPM Bind]
\label{def:rpm-bind}
The \textbf{bind} operation $(\rpmbind) : \RPM(A) \times (A \to \RPM(B)) \to \RPM(B)$ is:
\[
m \rpmbind f = \lambda \sigma. \begin{cases}
f(a)(\sigma') & \text{if } m(\sigma) = (\mathsf{Success}(a), \sigma') \\
(\mathsf{Failure}(o), \sigma') & \text{if } m(\sigma) = (\mathsf{Failure}(o), \sigma')
\end{cases}
\]
If the first computation succeeds, pass its result to $f$. If it fails, propagate the failure.
\end{definition}

\begin{definition}[RPM Multiplication]
\label{def:rpm-mult}
The \textbf{multiplication} $\mu : \RPM(\RPM(A)) \to \RPM(A)$ is:
\[
\mu(m) = m \rpmbind \id
\]
This flattens nested RPM computations.
\end{definition}

\section{Proof of Monad Laws}

\begin{theorem}[RPM Satisfies Monad Laws]
\label{thm:rpm-monad-laws}
The structure $(\RPM, \eta, \mu)$ satisfies the three monad laws.
\end{theorem}

\begin{proof}
\textbf{(M1) Left Unit Law:} $\rpmret(a) \rpmbind f = f(a)$

\begin{align}
(\rpmret(a) \rpmbind f)(\sigma) &= (\lambda \sigma'. (\mathsf{Success}(a), \sigma'))(\sigma) \text{ then } f \\
&= f(a)(\sigma)
\end{align}
Since $\rpmret(a)(\sigma) = (\mathsf{Success}(a), \sigma)$, we have $f(a)(\sigma)$.

\textbf{(M2) Right Unit Law:} $m \rpmbind \rpmret = m$

For any $m$:
\begin{align}
(m \rpmbind \rpmret)(\sigma) &= \begin{cases}
\rpmret(a)(\sigma') = (\mathsf{Success}(a), \sigma') & \text{if } m(\sigma) = (\mathsf{Success}(a), \sigma') \\
(\mathsf{Failure}(o), \sigma') & \text{if } m(\sigma) = (\mathsf{Failure}(o), \sigma')
\end{cases} \\
&= m(\sigma)
\end{align}

\textbf{(M3) Associativity:} $(m \rpmbind f) \rpmbind g = m \rpmbind (\lambda a. f(a) \rpmbind g)$

Let $m(\sigma) = (\mathsf{Success}(a), \sigma_1)$, $f(a)(\sigma_1) = (\mathsf{Success}(b), \sigma_2)$.
\begin{align}
\text{LHS: } ((m \rpmbind f) \rpmbind g)(\sigma) &= (f(a) \rpmbind g)(\sigma_1) = g(b)(\sigma_2) \\
\text{RHS: } (m \rpmbind (\lambda a. f(a) \rpmbind g))(\sigma) &= (f(a) \rpmbind g)(\sigma_1) = g(b)(\sigma_2)
\end{align}
Both sides equal. Failure cases propagate identically by definition of bind.
\end{proof}

\section{Prerequisite Checking Operations}

\begin{definition}[Check Acquisition]
\label{def:check-acq}
The operation $\mathsf{check} : \Acquisitions \to \RPM(\mathsf{Unit})$ verifies a prerequisite:
\[
\mathsf{check}(X) = \lambda \sigma. \begin{cases}
(\mathsf{Success}(()), \sigma) & \text{if } \sigma(X) = \mathtt{true} \\
(\mathsf{Failure}(X), \sigma) & \text{if } \sigma(X) = \mathtt{false}
\end{cases}
\]
\end{definition}

\begin{definition}[RPM Chain for Foundation $m$]
\label{def:rpm-chain}
The complete prerequisite chain for Foundation $m$ is:
\begin{align}
\mathsf{evalChain}_m &= \mathsf{check}(A0) \rpmbind \lambda \_. \\
&\quad\; \mathsf{check}(D_m) \rpmbind \lambda \_. \\
&\quad\; \mathsf{check}(W_m) \rpmbind \lambda \_. \\
&\quad\; \mathsf{check}(P_m) \rpmbind \lambda \_. \\
&\quad\; \rpmret(\mathsf{Outcome}_m)
\end{align}
\end{definition}

\begin{theorem}[Chain Evaluation Semantics]
\label{thm:chain-eval}
For an acquisition state $\sigma$:
\[
\mathsf{evalChain}_m(\sigma) = \begin{cases}
(\mathsf{Success}(\mathsf{Outcome}_m), \sigma) & \text{if all prerequisites satisfied} \\
(\mathsf{Failure}(X), \sigma) & \text{where } X \text{ is the first unsatisfied prerequisite}
\end{cases}
\]
\end{theorem}

\begin{proof}
By the definition of bind, the chain short-circuits at the first failure. If $\sigma(A0) = \mathtt{false}$, the result is $\mathsf{Failure}(A0)$. Otherwise, proceed to check $D_m$, and so on.
\end{proof}

\section{Dynamic Evaluation Rules}

\begin{definition}[RPM Evaluation Judgment]
\label{def:rpm-eval-judgment}
The evaluation judgment for RPM computations is:
\[
\langle E, \sigma \rangle \bigstep_\RPM v \mid \sigma'
\]
meaning: expression $E$ in state $\sigma$ evaluates to value $v$ with final state $\sigma'$.

For failures:
\[
\langle E, \sigma \rangle \bigstep_\RPM \Failed(o) \mid \sigma'
\]
\end{definition}

\begin{definition}[Big-Step Rules for RPM]
\label{def:rpm-bigstep}
\textbf{Return Rule:}
\[
\frac{}{\langle \rpmret(v), \sigma \rangle \bigstep_\RPM v \mid \sigma} \quad [\textsc{E-Return}]
\]

\textbf{Bind Success Rule:}
\[
\frac{\langle m, \sigma \rangle \bigstep_\RPM v \mid \sigma' \quad \langle f(v), \sigma' \rangle \bigstep_\RPM w \mid \sigma''}{\langle m \rpmbind f, \sigma \rangle \bigstep_\RPM w \mid \sigma''} \quad [\textsc{E-Bind-Succ}]
\]

\textbf{Bind Failure Rule:}
\[
\frac{\langle m, \sigma \rangle \bigstep_\RPM \Failed(o) \mid \sigma'}{\langle m \rpmbind f, \sigma \rangle \bigstep_\RPM \Failed(o) \mid \sigma'} \quad [\textsc{E-Bind-Fail}]
\]

\textbf{Check Success Rule:}
\[
\frac{\sigma(X) = \mathtt{true}}{\langle \mathsf{check}(X), \sigma \rangle \bigstep_\RPM () \mid \sigma} \quad [\textsc{E-Check-Succ}]
\]

\textbf{Check Failure Rule:}
\[
\frac{\sigma(X) = \mathtt{false}}{\langle \mathsf{check}(X), \sigma \rangle \bigstep_\RPM \Failed(X) \mid \sigma} \quad [\textsc{E-Check-Fail}]
\]
\end{definition}

\section{RPM Diagnostic Algorithm}

\begin{definition}[RPM Diagnostic Function]
\label{def:rpm-diag}
The diagnostic function identifies the first unsatisfied prerequisite:
\begin{align}
&\mathsf{diagnose} : \Foundations \times \AcqState \to \Acquisitions + \mathsf{Success} \\
&\mathsf{diagnose}(m, \sigma) = \\
&\quad \mathbf{if}\ \neg\sigma(A0)\ \mathbf{then}\ A0 \\
&\quad \mathbf{else\ if}\ \neg\sigma(D_m)\ \mathbf{then}\ D_m \\
&\quad \mathbf{else\ if}\ \neg\sigma(W_m)\ \mathbf{then}\ W_m \\
&\quad \mathbf{else\ if}\ \neg\sigma(P_m)\ \mathbf{then}\ P_m \\
&\quad \mathbf{else}\ \mathsf{Success}
\end{align}
\end{definition}

\begin{theorem}[RPM Termination]
\label{thm:rpm-term}
The RPM diagnostic algorithm terminates for any input in $O(1)$ time (constant number of checks).
\end{theorem}

\begin{proof}
The diagnostic performs at most 4 boolean checks: $A0, D_m, W_m, P_m$. Each check is $O(1)$. The total complexity is $O(4) = O(1)$.
\end{proof}

\section{Worked Examples: RPM Evaluation (v6.1)}

\begin{example}[Successful RPM Chain]
\label{ex:rpm-success}
\textbf{Goal:} Evaluate wealth acquisition ($F_6$) with all prerequisites satisfied.

\textbf{State:} $\sigma = \{A0 \mapsto \mathtt{true}, D_6 \mapsto \mathtt{true}, W_6 \mapsto \mathtt{true}, P_6 \mapsto \mathtt{true}\}$

\textbf{Evaluation:}
\begin{align}
&\langle \mathsf{evalChain}_6, \sigma \rangle \\
&= \langle \mathsf{check}(A0) \rpmbind \ldots, \sigma \rangle \\
&\bigstep_\RPM () \mid \sigma \quad [\text{check } A0 \text{ succeeds}] \\
&\bigstep_\RPM () \mid \sigma \quad [\text{check } D_6 \text{ succeeds}] \\
&\bigstep_\RPM () \mid \sigma \quad [\text{check } W_6 \text{ succeeds}] \\
&\bigstep_\RPM () \mid \sigma \quad [\text{check } P_6 \text{ succeeds}] \\
&\bigstep_\RPM \mathsf{Outcome}_6 \mid \sigma \quad [\text{return outcome}]
\end{align}

\textbf{Result:} $\mathsf{Success}(\mathsf{Outcome}_6)$ --- Wealth can manifest.
\end{example}

\begin{example}[Failed RPM Chain]
\label{ex:rpm-failure}
\textbf{Goal:} Evaluate wealth acquisition ($F_6$) with wisdom failure.

\textbf{State:} $\sigma = \{A0 \mapsto \mathtt{true}, D_6 \mapsto \mathtt{true}, W_6 \mapsto \mathtt{false}, P_6 \mapsto \mathtt{true}\}$

\textbf{Evaluation:}
\begin{align}
&\langle \mathsf{evalChain}_6, \sigma \rangle \\
&\bigstep_\RPM () \mid \sigma \quad [\text{check } A0 \text{ succeeds}] \\
&\bigstep_\RPM () \mid \sigma \quad [\text{check } D_6 \text{ succeeds}] \\
&\bigstep_\RPM \Failed(W_6) \mid \sigma \quad [\text{check } W_6 \text{ fails!}]
\end{align}

\textbf{Result:} $\mathsf{Failure}(W_6)$ --- Wisdom about wealth is missing.

\textbf{Diagnostic:} The person desires wealth and has power to act, but lacks correct mental models about money (perhaps believes ``money is evil'').

\textbf{Correction:} Repair $W_6$ through education, belief change, or therapy.
\end{example}

\begin{example}[RPM with Expression Evaluation]
\label{ex:rpm-expr}
\textbf{Expression:} $\RPM(A8^1_{6d})$ --- Evaluate esoteric truth in material foundation with RPM.

\textbf{State:} $\sigma = \{A0 \mapsto \mathtt{true}, D_6 \mapsto \mathtt{true}, W_6 \mapsto \mathtt{true}, P_6 \mapsto \mathtt{false}\}$

\textbf{Evaluation:}
\begin{enumerate}
\item Parse: $\RPM(A8^1_{6d})$
\item Extract Foundation: $m = 6$ (Material)
\item Run prerequisite chain:
   \begin{itemize}
   \item Check $A0$: $\checkmark$
   \item Check $D_6$: $\checkmark$
   \item Check $W_6$: $\checkmark$
   \item Check $P_6$: $\times$ \textbf{FAIL}
   \end{itemize}
\item Result: $\mathsf{Failure}(P_6)$
\end{enumerate}

\textbf{Interpretation:} The person has desire and wisdom for material success, but lacks power (action, resources, connections) to manifest.
\end{example}

%% --- NEW v7.0 CONTENT FOR RPM ---

\section{The Kleisli Category for RPM}
\label{sec:rpm-kleisli}

\vnew\ The RPM monad induces a Kleisli category that provides an alternative view of prerequisite-aware computations.

\begin{definition}[Kleisli Category $\Kleisli_\RPM$]
\label{def:kleisli-rpm}
The \textbf{Kleisli category} for the RPM monad is defined as:

\textbf{Objects:} $\Ob(\Kleisli_\RPM) = \Ob(\mathbf{Set})$ (all sets/types)

\textbf{Morphisms:} For objects $A, B$:
\[
\Hom_{\Kleisli}(A, B) = \Hom_{\mathbf{Set}}(A, \RPM(B))
\]
A Kleisli morphism $f : A \to_\Kleisli B$ is a function $f : A \to \RPM(B)$.

\textbf{Identity:} For each object $A$:
\[
\id_A^\Kleisli = \eta_A : A \to \RPM(A)
\]
where $\eta$ is the monad unit (return).

\textbf{Composition:} For $f : A \to_\Kleisli B$ and $g : B \to_\Kleisli C$:
\[
(g \circ_\Kleisli f)(a) = f(a) \rpmbind g = \mu_C(\RPM(g)(f(a)))
\]
\end{definition}

\begin{theorem}[Kleisli Category Axioms]
\label{thm:kleisli-axioms}
$\Kleisli_\RPM$ satisfies the category axioms:
\begin{enumerate}
  \item \textbf{Left Identity:} $\id_B^\Kleisli \circ_\Kleisli f = f$
  \item \textbf{Right Identity:} $f \circ_\Kleisli \id_A^\Kleisli = f$
  \item \textbf{Associativity:} $(h \circ_\Kleisli g) \circ_\Kleisli f = h \circ_\Kleisli (g \circ_\Kleisli f)$
\end{enumerate}
\end{theorem}

\begin{proof}
These follow directly from the monad laws for $\RPM$:

\textbf{Left Identity:}
\begin{align}
(\id_B^\Kleisli \circ_\Kleisli f)(a) &= f(a) \rpmbind \eta_B \\
&= f(a) \quad \text{(right unit law)}
\end{align}

\textbf{Right Identity:}
\begin{align}
(f \circ_\Kleisli \id_A^\Kleisli)(a) &= \eta_A(a) \rpmbind f \\
&= f(a) \quad \text{(left unit law)}
\end{align}

\textbf{Associativity:}
\begin{align}
((h \circ_\Kleisli g) \circ_\Kleisli f)(a) &= f(a) \rpmbind (\lambda b. g(b) \rpmbind h) \\
&= (f(a) \rpmbind g) \rpmbind h \quad \text{(associativity law)} \\
&= (h \circ_\Kleisli (g \circ_\Kleisli f))(a)
\end{align}
\end{proof}

\begin{definition}[Kleisli Lifting]
\label{def:kleisli-lift}
For a pure function $f : A \to B$, the \textbf{Kleisli lifting} is:
\[
f^\sharp : A \to_\Kleisli B = \eta_B \circ f
\]
Equivalently: $f^\sharp(a) = \rpmret(f(a))$
\end{definition}

\begin{proposition}[Lifting Preserves Composition]
\label{prop:lift-comp}
For pure functions $f : A \to B$ and $g : B \to C$:
\[
(g \circ f)^\sharp = g^\sharp \circ_\Kleisli f^\sharp
\]
\end{proposition}

\section{RPM as an Algebraic Effect}
\label{sec:rpm-effects}

\vnew\ The RPM monad can be understood as an algebraic effect with explicit effect operations and handlers.

\begin{definition}[RPM Effect Signature]
\label{def:rpm-effect-sig}
The \textbf{RPM effect signature} $\Sigma_\RPM$ consists of:

\textbf{Operations:}
\begin{align}
\mathsf{check} &: \Acquisitions \to \mathsf{Unit} \\
\mathsf{acquire} &: \Acquisitions \to \mathsf{Unit} \\
\mathsf{fail} &: \Acquisitions \to \bot
\end{align}

\textbf{Equations:}
\begin{align}
\mathsf{check}(p); \mathsf{check}(p) &= \mathsf{check}(p) \quad \text{(idempotence)} \\
\mathsf{acquire}(p); \mathsf{check}(p) &= \mathsf{acquire}(p) \quad \text{(acquire implies check)} \\
\mathsf{fail}(p); k &= \mathsf{fail}(p) \quad \text{(failure absorption)}
\end{align}
\end{definition}

\begin{definition}[RPM Effect Handler]
\label{def:rpm-handler}
An \textbf{RPM effect handler} interprets effect operations:
\begin{align}
\mathsf{handler}_\RPM &= \{ \\
&\quad \mathsf{return}(v) \mapsto \lambda\sigma. (\mathsf{Success}(v), \sigma) \\
&\quad \mathsf{check}(p, k) \mapsto \lambda\sigma.
  \begin{cases}
    k(())(\sigma) & \text{if } \sigma(p) \\
    (\mathsf{Failure}(p), \sigma) & \text{otherwise}
  \end{cases} \\
&\quad \mathsf{acquire}(p, k) \mapsto \lambda\sigma. k(())(\sigma[p \mapsto \mathsf{true}]) \\
&\quad \mathsf{fail}(p) \mapsto \lambda\sigma. (\mathsf{Failure}(p), \sigma) \\
&\}
\end{align}
where $k$ is the continuation.
\end{definition}

\section{Multi-Goal RPM Structures}
\label{sec:rpm-multigoal}

\vnew\ Real-world applications often require checking multiple prerequisites in parallel or pursuing multiple goals simultaneously.

\begin{definition}[Parallel RPM Composition]
\label{def:rpm-parallel}
For independent RPM computations $m_1 : \RPM(A)$ and $m_2 : \RPM(B)$:
\[
m_1 \pcomp m_2 : \RPM(A \times B)
\]
\[
(m_1 \pcomp m_2)(\sigma) =
\begin{cases}
(\mathsf{Success}(a, b), \sigma_2) & \text{if } m_1(\sigma) = (\mathsf{Success}(a), \sigma_1) \\
& \text{and } m_2(\sigma_1) = (\mathsf{Success}(b), \sigma_2) \\
(\mathsf{Failure}(p), \sigma') & \text{if either fails with } p
\end{cases}
\]
\end{definition}

\begin{theorem}[Parallel RPM is Commutative (for Independent Checks)]
\label{thm:rpm-parallel-comm}
If $m_1$ and $m_2$ check disjoint prerequisite sets, then:
\[
m_1 \pcomp m_2 \cong m_2 \pcomp m_1
\]
(isomorphic up to projection reordering)
\end{theorem}

\begin{definition}[RPM Choice]
\label{def:rpm-choice}
For alternative RPM computations:
\[
m_1 \oplus m_2 : \RPM(A)
\]
\[
(m_1 \oplus m_2)(\sigma) =
\begin{cases}
m_1(\sigma) & \text{if } m_1(\sigma) = (\mathsf{Success}(a), \sigma') \\
m_2(\sigma) & \text{if } m_1(\sigma) = (\mathsf{Failure}(\_), \_)
\end{cases}
\]
Try $m_1$; if it fails, try $m_2$.
\end{definition}

%% --- END NEW v7.0 RPM CONTENT ---

\chapter{Noetic Fractal Calculus: Complete Theory}
\label{ch:fractal-calculus}

\begin{tcolorbox}[colback=green!5!white,colframe=green!75!black,title=\vexp]
This chapter expands the v6.1 fractal calculus with complete type theory, eigenfractal analysis, convergence proofs, and categorical structure.
\end{tcolorbox}

\section{Fractals as Higher-Order Operators}

\begin{definition}[Noetic Fractal]
\label{def:noetic-fractal}
A \textbf{Noetic Fractal} $\Frac$ is a finite sequence of Noetic indices:
\[
\Frac = \langle k_1 : k_2 : \cdots : k_n \rangle
\]
where each $k_i \in \{0, 1, \ldots, 9\}$ and $n \geq 1$.
\end{definition}

\begin{definition}[Fractal Space]
\label{def:fractal-space}
The \textbf{Fractal space} is:
\[
\FracSet = \{\langle k_1 : \cdots : k_n \rangle : n \geq 1, k_i \in \{0, \ldots, 9\}\}
\]
\end{definition}

\begin{definition}[Fractal as Operator]
\label{def:fractal-operator}
Each fractal $\Frac = \langle k_1 : k_2 : \cdots : k_n \rangle$ defines an operator:
\[
\Frac : \Ideas \to \Ideas
\]
\[
\Frac(E) = \noe{k_1}(\noe{k_2}(\cdots\noe{k_n}(E)\cdots))
\]
Application is right-to-left (innermost first).
\end{definition}

\begin{definition}[Fractal Type]
\label{def:fractal-type}
In the TKS type system, fractals have type:
\[
\mathtt{Fractal} : \mathtt{List}^+(\mathbb{Z}_{10}) \to (\mathtt{Expr} \to \mathtt{Expr})
\]
where $\mathtt{List}^+$ denotes non-empty lists.
\end{definition}

\section{Fractal Composition}

\begin{definition}[Fractal Composition]
\label{def:fractal-comp}
For fractals $\Frac_1 = \langle a_1 : \cdots : a_m \rangle$ and $\Frac_2 = \langle b_1 : \cdots : b_n \rangle$:
\[
\Frac_1 \circ \Frac_2 = \langle a_1 : \cdots : a_m : b_1 : \cdots : b_n \rangle
\]
Concatenation of sequences.
\end{definition}

\begin{theorem}[Fractal Composition is Associative]
\label{thm:fractal-assoc}
For fractals $\Frac_1, \Frac_2, \Frac_3$:
\[
(\Frac_1 \circ \Frac_2) \circ \Frac_3 = \Frac_1 \circ (\Frac_2 \circ \Frac_3)
\]
\end{theorem}

\begin{proof}
List concatenation is associative.
\end{proof}

\begin{definition}[Identity Fractal]
\label{def:id-fractal}
The identity fractal is:
\[
\Frac_{\id} = \langle 0 \rangle
\]
\[
\Frac_{\id}(E) = \noe{0}(E) = E
\]
\end{definition}

\begin{theorem}[Fractal Monoid]
\label{thm:fractal-monoid}
The structure $(\FracSet, \circ, \Frac_{\id})$ forms a monoid.
\end{theorem}

\section{Fractal Simplification}

\begin{definition}[Dual Cancellation]
\label{def:dual-cancel}
Adjacent dual pairs simplify:
\begin{align}
\langle \cdots : 2 : 3 : \cdots \rangle &\Rightarrow \langle \cdots : 0 : \cdots \rangle \\
\langle \cdots : 3 : 2 : \cdots \rangle &\Rightarrow \langle \cdots : 0 : \cdots \rangle \\
\langle \cdots : 5 : 6 : \cdots \rangle &\Rightarrow \langle \cdots : 0 : \cdots \rangle \\
\langle \cdots : 6 : 5 : \cdots \rangle &\Rightarrow \langle \cdots : 0 : \cdots \rangle \\
\langle \cdots : 8 : 9 : \cdots \rangle &\Rightarrow \langle \cdots : 0 : \cdots \rangle \\
\langle \cdots : 9 : 8 : \cdots \rangle &\Rightarrow \langle \cdots : 0 : \cdots \rangle
\end{align}
\end{definition}

\begin{definition}[Identity Elimination]
\label{def:id-elim}
Identity elements can be removed:
\[
\langle \cdots : 0 : \cdots \rangle \Rightarrow \langle \cdots : \cdots \rangle
\]
(unless it's the only element)
\end{definition}

\begin{definition}[Canonical Form]
\label{def:canonical-form}
A fractal is in \textbf{canonical form} if:
\begin{enumerate}
\item No adjacent dual pairs exist.
\item No interior zeros exist (only $\langle 0 \rangle$ itself).
\item The sequence is minimal under these reductions.
\end{enumerate}
\end{definition}

\begin{theorem}[Canonical Form Existence]
\label{thm:canonical-exists}
Every fractal has a unique canonical form.
\end{theorem}

\begin{proof}
The simplification rules are confluent (any order of application yields the same result) and terminating (each rule reduces sequence length or removes a pair). By Newman's lemma, a unique normal form exists.
\end{proof}

\section{Fractal Semantics}

\begin{definition}[Fractal Denotation]
\label{def:fractal-denote}
The denotational semantics of a fractal is:
\[
\sem{\langle k_1 : k_2 : \cdots : k_n \rangle(E)} = \mathfrak{n}_{k_1}(\mathfrak{n}_{k_2}(\cdots\mathfrak{n}_{k_n}(\sem{E})\cdots))
\]
\end{definition}

\begin{theorem}[Compositional Fractal Semantics]
\label{thm:fractal-compositional}
\[
\sem{(\Frac_1 \circ \Frac_2)(E)} = \sem{\Frac_1}(\sem{\Frac_2(E)})
\]
\end{theorem}

\begin{proof}
By definition of fractal composition and the compositional nature of Noetic semantics.
\end{proof}

\section{Fractal Iteration and Convergence (v6.1)}

\begin{definition}[Iterated Fractal]
\label{def:iterated-fractal}
For fractal $\Frac$ and $n \geq 0$:
\[
\Frac^n = \underbrace{\Frac \circ \Frac \circ \cdots \circ \Frac}_{n \text{ times}}
\]
with $\Frac^0 = \Frac_{\id}$.
\end{definition}

\begin{definition}[Fractal Orbit]
\label{def:fractal-orbit}
The \textbf{orbit} of $E$ under $\Frac$ is:
\[
\mathcal{O}_\Frac(E) = \{E, \Frac(E), \Frac^2(E), \Frac^3(E), \ldots\}
\]
\end{definition}

\begin{definition}[Fractal Fixed Point]
\label{def:fractal-fixpoint}
$E^*$ is a \textbf{fixed point} of $\Frac$ if:
\[
\Frac(E^*) = E^*
\]
\end{definition}

\begin{definition}[Eigenfractal]
\label{def:eigenfractal}
$E$ is an \textbf{eigenstate} of $\Frac$ with eigenvalue $\lambda$ if:
\[
\Frac(E) = \lambda E
\]
Fixed points have $\lambda = 1$.
\end{definition}

\begin{definition}[Fractal Convergence]
\label{def:fractal-convergence}
$\Frac$ \textbf{converges} on $E$ if:
\[
\lim_{n \to \infty} \Frac^n(E) = E^*
\]
exists (in a suitable topology on $\Ideas$).
\end{definition}

\section{Fractal Classification}

\begin{definition}[Fractal Classes]
\label{def:fractal-classes}
Fractals are classified by their dynamic behavior:

\textbf{Stabilizing:} $\Frac$ is stabilizing if for all compatible $E$:
\[
\exists E^* : \lim_{n \to \infty} \Frac^n(E) = E^*
\]

\textbf{Amplifying:} $\Frac$ is amplifying if:
\[
\|\Frac^n(E)\| \to \infty \text{ as } n \to \infty
\]

\textbf{Oscillatory:} $\Frac$ is oscillatory if the orbit is periodic:
\[
\exists p > 0 : \Frac^p(E) = E
\]

\textbf{Neutral:} $\Frac$ is neutral if $\Frac(E) = E$ for all $E$ (identity).
\end{definition}

\begin{theorem}[MVR Stability]
\label{thm:mvr-stability}
The MVR fractal $\Frac_{\text{MVR}} = \langle 1 : 4 : 7 \rangle$ is stabilizing.
\end{theorem}

\begin{proof}
MVR applies Mind ($\noe{1}$), Vibration ($\noe{4}$), and Rhythm ($\noe{7}$). None of these are in dual pairs with each other. The combination:
\begin{itemize}
\item $\noe{7}$: Establishes periodic structure (bounded oscillation)
\item $\noe{4}$: Charges with energy (bounded amplitude)
\item $\noe{1}$: Applies awareness (stabilizing attention)
\end{itemize}
The composition creates a stable attractor pattern. Repeated application reinforces rather than amplifies.
\end{proof}

\begin{theorem}[Polarity Oscillation]
\label{thm:polarity-oscillation}
The polarity fractal $\Frac_{\text{pol}} = \langle 2 : 3 \rangle$ is oscillatory with period 2 (approximately).
\end{theorem}

\begin{proof}
$\langle 2 : 3 \rangle \circ \langle 2 : 3 \rangle = \langle 2 : 3 : 2 : 3 \rangle \Rightarrow \langle 0 : 0 \rangle \Rightarrow \langle 0 \rangle$.

After simplification, $\Frac_{\text{pol}}^2 \approx \Frac_{\id}$. The orbit returns to (approximately) the starting point.
\end{proof}

\section{Standard Fractal Families}

\begin{definition}[Standard Fractals]
\label{def:standard-fractals}
\begin{center}
\begin{tabular}{lll}
\toprule
\textbf{Name} & \textbf{Fractal} & \textbf{Class} \\
\midrule
MVR & $\langle 1 : 4 : 7 \rangle$ & Stabilizing \\
ACBE & $\langle 8 : 9 \rangle$ or cascade & Stabilizing \\
Polarity & $\langle 2 : 3 \rangle$ & Oscillatory \\
Gender & $\langle 5 : 6 \rangle$ & Oscillatory \\
Causation & $\langle 8 : 9 \rangle$ & Oscillatory \\
Deep Mind & $\langle 1 : 1 : 1 \rangle$ & Stabilizing \\
Resonance & $\langle 4 : 4 : 4 \rangle$ & Amplifying \\
Meta-Rhythm & $\langle 7 : 7 : 7 \rangle$ & Stabilizing \\
\bottomrule
\end{tabular}
\end{center}
\end{definition}

\section{Operational Semantics for Fractals (v6.1)}

\begin{definition}[Small-Step Fractal Evaluation]
\label{def:fractal-smallstep}
\[
\frac{}{\langle k \rangle(E) \smallstep \noe{k}(E)} \quad [\textsc{S-Frac-Single}]
\]

\[
\frac{}{\langle k_1 : k_2 : \cdots : k_n \rangle(E) \smallstep \noe{k_1}(\langle k_2 : \cdots : k_n \rangle(E))} \quad [\textsc{S-Frac-Step}]
\]
\end{definition}

\begin{definition}[Big-Step Fractal Evaluation]
\label{def:fractal-bigstep}
\[
\frac{E \bigstep v}{\langle k \rangle(E) \bigstep \mathfrak{n}_k(v)} \quad [\textsc{E-Frac-Single}]
\]

\[
\frac{E \bigstep v_n \quad \langle k_1 : \cdots : k_{n-1} \rangle(v_n) \bigstep v_1}{\langle k_1 : \cdots : k_n \rangle(E) \bigstep v_1} \quad [\textsc{E-Frac-Multi}]
\]
\end{definition}

\section{Worked Examples: Fractal Evaluation (v6.1)}

\begin{example}[MVR Application]
\label{ex:mvr-application}
\textbf{Expression:} $\langle 1 : 4 : 7 \rangle(D2_{3d})$

\textbf{Evaluation:}
\begin{align}
&\langle 1 : 4 : 7 \rangle(D2_{3d}) \\
&\smallstep \noe{1}(\langle 4 : 7 \rangle(D2_{3d})) \\
&\smallstep \noe{1}(\noe{4}(\langle 7 \rangle(D2_{3d}))) \\
&\smallstep \noe{1}(\noe{4}(\noe{7}(D2_{3d}))) \\
&\bigstep \noe{1}(\noe{4}(\text{Rhythmic Physical Positive})) \\
&\bigstep \noe{1}(\text{Vibrating Rhythmic Physical Positive}) \\
&\bigstep \text{Aware Vibrating Rhythmic Physical Positive}
\end{align}

\textbf{Semantic interpretation:} A conscious, vibratory, rhythmic pattern of physical health (for Life foundation).
\end{example}

\begin{example}[Polarity Cancellation]
\label{ex:polarity-cancel}
\textbf{Expression:} $\langle 2 : 3 : 2 : 3 \rangle(A8)$

\textbf{Simplification:}
\[
\langle 2 : 3 : 2 : 3 \rangle \Rightarrow \langle 0 : 0 \rangle \Rightarrow \langle 0 \rangle
\]

\textbf{Evaluation:}
\[
\langle 0 \rangle(A8) = \noe{0}(A8) = A8
\]

\textbf{Result:} The polarity oscillations cancel, returning to the original expression.
\end{example}

\begin{example}[Deep Meditation Fractal]
\label{ex:deep-meditation}
\textbf{Expression:} $\langle 1 : 1 : 1 : 4 : 4 : 4 : 7 : 7 : 7 \rangle(A1_{1a})$

\textbf{Interpretation:} Triple-strength MVR for deep spiritual practice.

\textbf{Evaluation:}
\begin{align}
&= (\noe{1})^3 \circ (\noe{4})^3 \circ (\noe{7})^3 (A1_{1a}) \\
&= \noe{1}(\noe{1}(\noe{1}(\noe{4}(\noe{4}(\noe{4}(\noe{7}(\noe{7}(\noe{7}(A1_{1a})))))))))
\end{align}

\textbf{Result:} Deeply stabilized spiritual awareness with strong rhythmic and vibratory components.
\end{example}

%% --- NEW v7.0 CONTENT FOR FRACTALS ---

\section{Fractal Type Theory}
\label{sec:fractal-types}

\vnew\ Fractals have a rich type structure that can be captured in an indexed type system.

\begin{definition}[Fractal Type Constructor]
\label{def:fractal-type-ctor}
The fractal type constructor is:
\[
\mathsf{Frac} : \mathsf{List}^+(\mathbb{Z}_{10}) \to \mathsf{Type}
\]
For a noetic sequence $\bar{k} = \langle k_1 : k_2 : \cdots : k_n \rangle$:
\[
\mathsf{Frac}[\bar{k}] : \mathsf{Domain} \to \mathsf{Domain}
\]
\end{definition}

\begin{definition}[Fractal Kinding Rules]
\label{def:fractal-kinding}
\[
\frac{k_i \in \{0, 1, \ldots, 9\} \quad \forall i \in \{1, \ldots, n\}}
{\mathsf{Frac}[\langle k_1 : \cdots : k_n \rangle] : \mathsf{Type}} \quad [\textsc{K-Frac}]
\]
\end{definition}

\begin{definition}[Fractal Type Operations]
\label{def:fractal-type-ops}
Type-level operations on fractal types:

\textbf{Concatenation:}
\[
\mathsf{Frac}[\bar{k}_1] \cdot \mathsf{Frac}[\bar{k}_2] = \mathsf{Frac}[\bar{k}_1 \mathbin{+\!\!+} \bar{k}_2]
\]

\textbf{Length:}
\[
|\mathsf{Frac}[\langle k_1 : \cdots : k_n \rangle]| = n
\]

\textbf{Projection:}
\[
\pi_i(\mathsf{Frac}[\langle k_1 : \cdots : k_n \rangle]) = k_i
\]
\end{definition}

\section{Eigenfractal Analysis}
\label{sec:eigenfractals}

\vnew\ Eigenfractals are the fixed points and eigenstructures of fractal operators.

\begin{definition}[Fractal Eigenvalue]
\label{def:fractal-eigenvalue}
For fractal $\Frac$ and Idea $E$, $\lambda$ is an \textbf{eigenvalue} if:
\[
\Frac(E) = \lambda \cdot E
\]
where $\cdot$ denotes scaling in the Idea space (when defined).
\end{definition}

\begin{definition}[Eigenvalue Spectrum]
\label{def:eigenvalue-spectrum}
The \textbf{spectrum} of fractal $\Frac$ is:
\[
\Spec(\Frac) = \{\lambda : \exists E \neq \varnothing, \Frac(E) = \lambda \cdot E\}
\]
\end{definition}

\begin{theorem}[Fixed Point Existence for Stabilizing Fractals]
\label{thm:stabilizing-fixpoint}
Every stabilizing fractal $\Frac$ has at least one fixed point in any non-empty closed invariant subspace of $\Ideas$.
\end{theorem}

\begin{proof}
Let $S \subseteq \Ideas$ be a non-empty closed invariant subspace (i.e., $\Frac(S) \subseteq S$).

Since $\Frac$ is stabilizing, for any $E_0 \in S$, the sequence:
\[
E_0, \Frac(E_0), \Frac^2(E_0), \ldots
\]
converges to some limit $E^* \in S$ (by closedness).

By continuity of $\Frac$ (which follows from compositionality of Noetic operators):
\[
\Frac(E^*) = \Frac(\lim_{n \to \infty} \Frac^n(E_0)) = \lim_{n \to \infty} \Frac^{n+1}(E_0) = E^*
\]
Thus $E^*$ is a fixed point.
\end{proof}

\begin{definition}[Eigenfractal Decomposition]
\label{def:eigen-decomposition}
For a fractal $\Frac$ with spectrum $\Spec(\Frac) = \{\lambda_1, \ldots, \lambda_m\}$, the \textbf{eigenspace decomposition} is:
\[
\Ideas = \bigoplus_{i=1}^{m} V_{\lambda_i}
\]
where $V_{\lambda_i} = \{E : \Frac(E) = \lambda_i \cdot E\}$ is the eigenspace for $\lambda_i$.
\end{definition}

\section{Convergence Analysis}
\label{sec:fractal-convergence}

\vnew\ Rigorous analysis of fractal iteration convergence.

\begin{definition}[Fractal Norm]
\label{def:fractal-norm}
For an Idea space equipped with a norm $\|\cdot\|$, the \textbf{operator norm} of fractal $\Frac$ is:
\[
\|\Frac\|_{op} = \sup_{E \neq \varnothing} \frac{\|\Frac(E)\|}{\|E\|}
\]
\end{definition}

\begin{theorem}[Contraction Convergence]
\label{thm:contraction-conv}
If $\|\Frac\|_{op} < 1$ (contraction), then for any $E_0 \in \Ideas$:
\begin{enumerate}
  \item The sequence $\Frac^n(E_0)$ converges to a unique fixed point $E^*$
  \item The convergence rate is geometric: $\|\Frac^n(E_0) - E^*\| \leq \|\Frac\|_{op}^n \|E_0 - E^*\|$
\end{enumerate}
\end{theorem}

\begin{proof}
By the Banach fixed-point theorem applied to the complete metric space $(\Ideas, d)$ where $d(E_1, E_2) = \|E_1 - E_2\|$.

Since $\|\Frac\|_{op} < 1$, $\Frac$ is a contraction mapping. Banach's theorem guarantees existence and uniqueness of the fixed point, with the stated convergence rate.
\end{proof}

\begin{theorem}[MVR Convergence Rate]
\label{thm:mvr-convergence}
The MVR fractal $\Frac_{MVR} = \langle 1:4:7 \rangle$ converges in at most $O(\log(1/\epsilon))$ iterations to within $\epsilon$ of its fixed point.
\end{theorem}

\section{Fractal Categories}
\label{sec:fractal-categories}

\vnew\ The space of fractals has categorical structure.

\begin{definition}[Category of Fractals]
\label{def:frac-category}
The category $\mathbf{Frac}$ has:

\textbf{Objects:} Fractal types $\mathsf{Frac}[\bar{k}]$

\textbf{Morphisms:} Fractal homomorphisms preserving structure
\[
\Hom_\mathbf{Frac}(\mathsf{Frac}[\bar{k}_1], \mathsf{Frac}[\bar{k}_2])
\]

\textbf{Identity:} $\id_{\mathsf{Frac}[\bar{k}]} = \mathsf{Frac}[\langle 0 \rangle]$ composed appropriately

\textbf{Composition:} Fractal composition $\circ$
\end{definition}

\begin{proposition}[Frac is a Monoid Category]
\label{prop:frac-monoid}
$(\mathbf{Frac}, \circ, \mathsf{Frac}[\langle 0 \rangle])$ forms a strict monoidal category.
\end{proposition}

%% --- END NEW v7.0 FRACTAL CONTENT ---

%% ============================================================================
%% PART V: COMPLETE DYNAMIC SEMANTICS (NEW IN v7.0)
%% ============================================================================
\part{Complete Dynamic Semantics}

\chapter{Evaluation Environments and Contexts}
\label{ch:eval-env}

\begin{tcolorbox}[colback=blue!5!white,colframe=blue!75!black,title=\vnew]
This chapter provides the formal foundation for TKS evaluation, defining environments, stores, and evaluation contexts.
\end{tcolorbox}

\section{TKS Evaluation Environment}

\begin{definition}[Complete TKS Environment]
\label{def:tks-environment}
A \textbf{TKS evaluation environment} is a tuple:
\[
\mathcal{E} = (\rho, \sigma, \alpha, \kappa)
\]
where:
\begin{itemize}
  \item $\rho : \mathrm{Var} \to \mathrm{Val}$ --- variable bindings
  \item $\sigma : \mathrm{Loc} \to \mathrm{Val}$ --- mutable store
  \item $\alpha : \Acquisitions \to \mathbb{B}$ --- acquisition state
  \item $\kappa : \FracSet^*$ --- fractal context stack
\end{itemize}
\end{definition}

\begin{definition}[Environment Operations]
\label{def:env-ops}
\textbf{Variable Lookup:}
\[
\mathcal{E}(x) = \rho(x) \quad \text{where } \mathcal{E} = (\rho, \_, \_, \_)
\]

\textbf{Variable Extension:}
\[
\mathcal{E}[x \mapsto v] = (\rho[x \mapsto v], \sigma, \alpha, \kappa)
\]

\textbf{Acquisition Check:}
\[
\mathcal{E} \models p \iff \alpha(p) = \mathsf{true}
\]

\textbf{Acquisition Update:}
\[
\mathcal{E}[p \leftarrow \mathsf{true}] = (\rho, \sigma, \alpha[p \mapsto \mathsf{true}], \kappa)
\]

\textbf{Fractal Push:}
\[
\mathcal{E}.\mathsf{push}(\Frac) = (\rho, \sigma, \alpha, \Frac :: \kappa)
\]
\end{definition}

\section{Evaluation Contexts}

\begin{definition}[Evaluation Context Grammar]
\label{def:eval-context}
Evaluation contexts $\mathbb{E}$ specify where reduction occurs:
\begin{align}
\mathbb{E} ::=&\; [\cdot] \\
    &\mid \mathbb{E} + e \mid v + \mathbb{E} \\
    &\mid \mathbb{E} \times e \mid v \times \mathbb{E} \\
    &\mid \noe{k}(\mathbb{E}) \\
    &\mid \langle \bar{k} \rangle(\mathbb{E}) \\
    &\mid \mathbb{E} \rpmbind f \mid \rpmret\; \mathbb{E} \\
    &\mid \mathsf{check}(\mathbb{E}) \mid \mathsf{acquire}(\mathbb{E}) \\
    &\mid \mathbb{E}\; e \mid v\; \mathbb{E} \\
    &\mid (\mathbb{E}, e) \mid (v, \mathbb{E}) \\
    &\mid \mathsf{fst}\; \mathbb{E} \mid \mathsf{snd}\; \mathbb{E}
\end{align}
\end{definition}

\begin{definition}[Context Application]
\label{def:context-app}
For context $\mathbb{E}$ and expression $e$, $\mathbb{E}[e]$ denotes substitution of $e$ for the hole $[\cdot]$.
\end{definition}

\chapter{Big-Step Semantics for All Constructs}
\label{ch:bigstep}

\begin{tcolorbox}[colback=blue!5!white,colframe=blue!75!black,title=\vnew]
Complete big-step (natural) semantics for every TKS construct.
\end{tcolorbox}

\section{Element Evaluation}

\begin{definition}[Element Big-Step Rules]
\label{def:element-bigstep}

\textbf{World-Mode Element:}
\[
\frac{W \in \{A, B, C, D\} \quad n \in \{1, \ldots, 10\}}
{\langle Wn, \mathcal{E} \rangle \bigstep (Wn, \mathcal{E})} \quad [\textsc{E-Element}]
\]

\textbf{Foundation Reference:}
\[
\frac{m \in \{1, \ldots, 7\} \quad w \in \{a, b, c, d\}}
{\langle F_{mw}, \mathcal{E} \rangle \bigstep (F_{mw}, \mathcal{E})} \quad [\textsc{E-Foundation}]
\]
\end{definition}

\section{Noetic Evaluation}

\begin{definition}[Noetic Big-Step Rules]
\label{def:noetic-bigstep}

\textbf{Noetic Application:}
\[
\frac{\langle e, \mathcal{E} \rangle \bigstep (v, \mathcal{E}')}
{\langle \noe{k}(e), \mathcal{E} \rangle \bigstep (\mathfrak{n}_k(v), \mathcal{E}')} \quad [\textsc{E-Noe}]
\]

\textbf{Noetic Composition:}
\[
\frac{\langle e, \mathcal{E} \rangle \bigstep (v, \mathcal{E}') \quad \langle \noe{j}(\noe{k}(v)), \mathcal{E}' \rangle \bigstep (v', \mathcal{E}'')}
{\langle (\noe{j} \circ \noe{k})(e), \mathcal{E} \rangle \bigstep (v', \mathcal{E}'')} \quad [\textsc{E-Noe-Comp}]
\]

\textbf{Identity Noetic:}
\[
\frac{\langle e, \mathcal{E} \rangle \bigstep (v, \mathcal{E}')}
{\langle \noe{0}(e), \mathcal{E} \rangle \bigstep (v, \mathcal{E}')} \quad [\textsc{E-Noe-Id}]
\]
\end{definition}

\section{Tootra Arithmetic Evaluation}

\begin{definition}[Tootra Arithmetic Big-Step Rules]
\label{def:tootra-bigstep}

\textbf{Tootra-Addition:}
\[
\frac{\langle e_1, \mathcal{E} \rangle \bigstep (v_1, \mathcal{E}') \quad \langle e_2, \mathcal{E}' \rangle \bigstep (v_2, \mathcal{E}'')}
{\langle e_1 \Tadd e_2, \mathcal{E} \rangle \bigstep (v_1 \uplus v_2, \mathcal{E}'')} \quad [\textsc{E-TAdd}]
\]

\textbf{Tootra-Subtraction:}
\[
\frac{\langle e_1, \mathcal{E} \rangle \bigstep (v_1, \mathcal{E}') \quad \langle e_2, \mathcal{E}' \rangle \bigstep (v_2, \mathcal{E}'')}
{\langle e_1 \Tsub e_2, \mathcal{E} \rangle \bigstep (v_1 \setminus_m v_2, \mathcal{E}'')} \quad [\textsc{E-TSub}]
\]

\textbf{Tootra-Multiplication:}
\[
\frac{\langle e_1, \mathcal{E} \rangle \bigstep (v_1, \mathcal{E}') \quad \langle e_2, \mathcal{E}' \rangle \bigstep (v_2, \mathcal{E}'')}
{\langle e_1 \Tmul e_2, \mathcal{E} \rangle \bigstep (v_1 \odot v_2, \mathcal{E}'')} \quad [\textsc{E-TMul}]
\]

\textbf{Tootra-Division:}
\[
\frac{\langle e_1, \mathcal{E} \rangle \bigstep (v_1, \mathcal{E}') \quad \langle e_2, \mathcal{E}' \rangle \bigstep (v_2, \mathcal{E}'') \quad v_2 \neq \varnothing}
{\langle e_1 \Tdiv e_2, \mathcal{E} \rangle \bigstep (v_1 \oslash v_2, \mathcal{E}'')} \quad [\textsc{E-TDiv}]
\]
\end{definition}

\section{Fractal Evaluation}

\begin{definition}[Fractal Big-Step Rules]
\label{def:fractal-bigstep-complete}

\textbf{Single-Element Fractal:}
\[
\frac{\langle e, \mathcal{E} \rangle \bigstep (v, \mathcal{E}')}
{\langle \langle k \rangle(e), \mathcal{E} \rangle \bigstep (\mathfrak{n}_k(v), \mathcal{E}')} \quad [\textsc{E-Frac-1}]
\]

\textbf{Multi-Element Fractal:}
\[
\frac{\langle \langle k_2 : \cdots : k_n \rangle(e), \mathcal{E} \rangle \bigstep (v', \mathcal{E}') \quad \mathfrak{n}_{k_1}(v') = v''}
{\langle \langle k_1 : k_2 : \cdots : k_n \rangle(e), \mathcal{E} \rangle \bigstep (v'', \mathcal{E}')} \quad [\textsc{E-Frac-N}]
\]

\textbf{Identity Fractal:}
\[
\frac{\langle e, \mathcal{E} \rangle \bigstep (v, \mathcal{E}')}
{\langle \langle 0 \rangle(e), \mathcal{E} \rangle \bigstep (v, \mathcal{E}')} \quad [\textsc{E-Frac-Id}]
\]

\textbf{Fractal Composition:}
\[
\frac{\langle \Frac_2(e), \mathcal{E} \rangle \bigstep (v', \mathcal{E}') \quad \langle \Frac_1(v'), \mathcal{E}' \rangle \bigstep (v'', \mathcal{E}'')}
{\langle (\Frac_1 \circ \Frac_2)(e), \mathcal{E} \rangle \bigstep (v'', \mathcal{E}'')} \quad [\textsc{E-Frac-Comp}]
\]
\end{definition}

\section{RPM Evaluation}

\textit{[NOTE: RPM evaluation rules from v6.1 preserved, see Chapter 7]}

\section{ACBE Functor Evaluation}

\begin{definition}[ACBE Big-Step Rules]
\label{def:acbe-bigstep}

\textbf{ACBE Application:}
\[
\frac{\langle e, \mathcal{E} \rangle \bigstep (An, \mathcal{E}')}
{\langle \mathrm{ACBE}(e), \mathcal{E} \rangle \bigstep (Dn, \mathcal{E}')} \quad [\textsc{E-ACBE}]
\]

\textbf{ACBE on Non-Spiritual:}
\[
\frac{\langle e, \mathcal{E} \rangle \bigstep (Wn, \mathcal{E}') \quad W \neq A}
{\langle \mathrm{ACBE}(e), \mathcal{E} \rangle \bigstep (\mathsf{error}, \mathcal{E}')} \quad [\textsc{E-ACBE-Err}]
\]

\textbf{Cascade ACBE:}
\[
\frac{\langle e, \mathcal{E} \rangle \bigstep (An, \mathcal{E}_0) \quad (An \to Bn \to Cn \to Dn)}
{\langle \mathrm{ACBE}_{cascade}(e), \mathcal{E} \rangle \bigstep (Dn, \mathcal{E}_0)} \quad [\textsc{E-ACBE-Casc}]
\]
\end{definition}

%% [CONTINUE WITH SMALL-STEP SEMANTICS, TYPE PROOFS, CONCURRENCY, QUANTUM...]

%% ============================================================================
%% PLACEHOLDER FOR REMAINING CHAPTERS
%% ============================================================================

\chapter{Small-Step Semantics}
\label{ch:smallstep}

\begin{tcolorbox}[colback=blue!5!white,colframe=blue!75!black,title=\vnew]
Complete small-step (structural operational) semantics for TKS, corresponding to the big-step rules in Chapter~\ref{ch:bigstep}.
\end{tcolorbox}

\section{Values and Redexes}

\begin{definition}[TKS Values]
\label{def:tks-values}
Values are fully evaluated expressions:
\begin{align}
v ::=&\; Wn \quad \text{(Element values, } W \in \{A,B,C,D\}, n \in \{1,\ldots,10\}) \\
    &\mid F_{mw} \quad \text{(Foundation values)} \\
    &\mid \lambda x:\tau. e \quad \text{(Function values)} \\
    &\mid (v_1, v_2) \quad \text{(Pair values)} \\
    &\mid \mathsf{inl}\; v \mid \mathsf{inr}\; v \quad \text{(Sum values)} \\
    &\mid \rpmret\; v \quad \text{(RPM success values)} \\
    &\mid \Failed(p) \quad \text{(RPM failure values)} \\
    &\mid n \quad \text{(Integer literals)} \\
    &\mid \mathsf{true} \mid \mathsf{false} \quad \text{(Boolean values)} \\
    &\mid () \quad \text{(Unit value)}
\end{align}
\end{definition}

\begin{definition}[Redexes]
\label{def:redexes}
A \textbf{redex} (reducible expression) is an expression that can take a step:
\begin{itemize}
  \item $(\lambda x:\tau. e)\; v$ --- beta redex
  \item $\noe{k}(v)$ --- noetic redex
  \item $\langle \bar{k} \rangle(v)$ --- fractal redex
  \item $v_1 \Tadd v_2$ --- arithmetic redex
  \item $(\rpmret\; v) \rpmbind f$ --- monadic redex
  \item $\mathsf{check}(p)$ --- prerequisite check redex
  \item $\mathsf{acquire}(p)$ --- acquisition redex
\end{itemize}
\end{definition}

\section{Core Reduction Rules}

\begin{definition}[Small-Step Reduction $\smallstep$]
\label{def:smallstep-rules}

\textbf{Beta Reduction:}
\[
\frac{}
{(\lambda x:\tau. e)\; v \smallstep e[v/x]} \quad [\textsc{S-Beta}]
\]

\textbf{Let Reduction:}
\[
\frac{}
{\mathsf{let}\; x = v \;\mathsf{in}\; e \smallstep e[v/x]} \quad [\textsc{S-Let}]
\]

\textbf{Conditional - True:}
\[
\frac{}
{\mathsf{if}\; \mathsf{true} \;\mathsf{then}\; e_1 \;\mathsf{else}\; e_2 \smallstep e_1} \quad [\textsc{S-IfT}]
\]

\textbf{Conditional - False:}
\[
\frac{}
{\mathsf{if}\; \mathsf{false} \;\mathsf{then}\; e_1 \;\mathsf{else}\; e_2 \smallstep e_2} \quad [\textsc{S-IfF}]
\]

\textbf{Projection:}
\[
\frac{}
{\mathsf{fst}\; (v_1, v_2) \smallstep v_1} \quad [\textsc{S-Fst}]
\qquad
\frac{}
{\mathsf{snd}\; (v_1, v_2) \smallstep v_2} \quad [\textsc{S-Snd}]
\]
\end{definition}

\section{Noetic Small-Step Rules}

\begin{definition}[Noetic Reduction]
\label{def:noetic-smallstep}

\textbf{Noetic Application:}
\[
\frac{}
{\noe{k}(v) \smallstep \mathfrak{n}_k(v)} \quad [\textsc{S-Noe}]
\]

\textbf{Noetic Congruence:}
\[
\frac{e \smallstep e'}
{\noe{k}(e) \smallstep \noe{k}(e')} \quad [\textsc{S-Noe-Cong}]
\]

\textbf{Identity Noetic:}
\[
\frac{}
{\noe{0}(v) \smallstep v} \quad [\textsc{S-Noe-Id}]
\]

\textbf{Noetic Composition:}
\[
\frac{}
{(\noe{j} \circ \noe{k})(v) \smallstep \noe{j}(\noe{k}(v))} \quad [\textsc{S-Noe-Comp}]
\]
\end{definition}

\section{Fractal Small-Step Rules}

\begin{definition}[Fractal Reduction]
\label{def:fractal-smallstep}

\textbf{Single-Element Fractal:}
\[
\frac{}
{\langle k \rangle(v) \smallstep \noe{k}(v)} \quad [\textsc{S-Frac-1}]
\]

\textbf{Multi-Element Fractal Expansion:}
\[
\frac{}
{\langle k_1 : k_2 : \cdots : k_n \rangle(v) \smallstep \noe{k_1}(\langle k_2 : \cdots : k_n \rangle(v))} \quad [\textsc{S-Frac-Exp}]
\]

\textbf{Fractal Congruence:}
\[
\frac{e \smallstep e'}
{\langle \bar{k} \rangle(e) \smallstep \langle \bar{k} \rangle(e')} \quad [\textsc{S-Frac-Cong}]
\]

\textbf{Identity Fractal:}
\[
\frac{}
{\langle 0 \rangle(v) \smallstep v} \quad [\textsc{S-Frac-Id}]
\]

\textbf{Fractal Composition:}
\[
\frac{}
{(\Frac_1 \circ \Frac_2)(v) \smallstep \Frac_1(\Frac_2(v))} \quad [\textsc{S-Frac-Comp}]
\]
\end{definition}

\section{Tootra Arithmetic Small-Step Rules}

\begin{definition}[Tootra Arithmetic Reduction]
\label{def:tootra-smallstep}

\textbf{Tootra Addition:}
\[
\frac{}
{v_1 \Tadd v_2 \smallstep v_1 \uplus v_2} \quad [\textsc{S-TAdd}]
\]

\textbf{Tootra Subtraction:}
\[
\frac{}
{v_1 \Tsub v_2 \smallstep v_1 \setminus_m v_2} \quad [\textsc{S-TSub}]
\]

\textbf{Tootra Multiplication:}
\[
\frac{}
{v_1 \Tmul v_2 \smallstep v_1 \odot v_2} \quad [\textsc{S-TMul}]
\]

\textbf{Tootra Division:}
\[
\frac{v_2 \neq \varnothing}
{v_1 \Tdiv v_2 \smallstep v_1 \oslash v_2} \quad [\textsc{S-TDiv}]
\]

\textbf{Left Congruence:}
\[
\frac{e_1 \smallstep e_1'}
{e_1 \Tadd e_2 \smallstep e_1' \Tadd e_2} \quad [\textsc{S-TAdd-L}]
\]

\textbf{Right Congruence:}
\[
\frac{e_2 \smallstep e_2'}
{v_1 \Tadd e_2 \smallstep v_1 \Tadd e_2'} \quad [\textsc{S-TAdd-R}]
\]
\end{definition}

\section{RPM Small-Step Rules}

\begin{definition}[RPM Reduction with State]
\label{def:rpm-smallstep}
RPM reductions carry state $\alpha$ (acquisition state):

\textbf{Return-Bind:}
\[
\frac{}
{(\rpmret\; v, \alpha) \smallstep_\RPM (f(v), \alpha) \text{ in } (\rpmret\; v) \rpmbind f} \quad [\textsc{S-Ret-Bind}]
\]

\textbf{Fail-Bind:}
\[
\frac{}
{(\Failed(p), \alpha) \rpmbind f \smallstep (\Failed(p), \alpha)} \quad [\textsc{S-Fail-Bind}]
\]

\textbf{Check Success:}
\[
\frac{p \in \alpha}
{(\mathsf{check}(p), \alpha) \smallstep (\rpmret\; (), \alpha)} \quad [\textsc{S-Check-Ok}]
\]

\textbf{Check Failure:}
\[
\frac{p \notin \alpha}
{(\mathsf{check}(p), \alpha) \smallstep (\Failed(p), \alpha)} \quad [\textsc{S-Check-Fail}]
\]

\textbf{Acquire:}
\[
\frac{}
{(\mathsf{acquire}(p), \alpha) \smallstep (\rpmret\; (), \alpha \cup \{p\})} \quad [\textsc{S-Acquire}]
\]

\textbf{Bind Congruence:}
\[
\frac{(m, \alpha) \smallstep (m', \alpha')}
{(m \rpmbind f, \alpha) \smallstep (m' \rpmbind f, \alpha')} \quad [\textsc{S-Bind-Cong}]
\]
\end{definition}

\section{ACBE Small-Step Rules}

\begin{definition}[ACBE Reduction]
\label{def:acbe-smallstep}

\textbf{ACBE Direct:}
\[
\frac{}
{\mathrm{ACBE}(An) \smallstep Dn} \quad [\textsc{S-ACBE}]
\]

\textbf{ACBE Cascade Step:}
\[
\frac{}
{\mathrm{ACBE}_{step}(An) \smallstep Bn} \quad [\textsc{S-ACBE-A}]
\]
\[
\frac{}
{\mathrm{ACBE}_{step}(Bn) \smallstep Cn} \quad [\textsc{S-ACBE-B}]
\]
\[
\frac{}
{\mathrm{ACBE}_{step}(Cn) \smallstep Dn} \quad [\textsc{S-ACBE-C}]
\]

\textbf{ACBE Congruence:}
\[
\frac{e \smallstep e'}
{\mathrm{ACBE}(e) \smallstep \mathrm{ACBE}(e')} \quad [\textsc{S-ACBE-Cong}]
\]
\end{definition}

\section{Context Rules}

\begin{definition}[Contextual Reduction]
\label{def:context-reduction}
For evaluation context $\mathbb{E}$ (Definition~\ref{def:eval-context}):
\[
\frac{e \smallstep e'}
{\mathbb{E}[e] \smallstep \mathbb{E}[e']} \quad [\textsc{S-Context}]
\]
\end{definition}

\section{Equivalence with Big-Step Semantics}

\begin{theorem}[Small-Step/Big-Step Equivalence]
\label{thm:step-equivalence}
For all expressions $e$, environments $\mathcal{E}$, values $v$, and final environments $\mathcal{E}'$:
\[
\langle e, \mathcal{E} \rangle \bigstep (v, \mathcal{E}') \iff \langle e, \mathcal{E} \rangle \smallstep^* (v, \mathcal{E}')
\]
where $\smallstep^*$ is the reflexive-transitive closure of $\smallstep$.
\end{theorem}

\begin{proof}
By mutual induction on big-step derivations (soundness) and small-step reduction sequences (completeness).

\textbf{Soundness ($\bigstep \Rightarrow \smallstep^*$):}
By induction on the big-step derivation. Each big-step rule corresponds to a sequence of small steps.

\textbf{Completeness ($\smallstep^* \Rightarrow \bigstep$):}
By induction on the length of the reduction sequence, using determinism of small-step evaluation.
\end{proof}

\begin{theorem}[Determinism]
\label{thm:determinism}
TKS small-step evaluation is deterministic:
\[
\text{If } e \smallstep e_1 \text{ and } e \smallstep e_2, \text{ then } e_1 = e_2
\]
\end{theorem}

\begin{proof}
By structural induction on expressions, observing that evaluation contexts uniquely determine the redex position and each redex has a unique contractum.
\end{proof}

\begin{theorem}[Confluence]
\label{thm:confluence}
TKS reduction is confluent (Church-Rosser):
\[
\text{If } e \smallstep^* e_1 \text{ and } e \smallstep^* e_2, \text{ then } \exists e'. e_1 \smallstep^* e' \text{ and } e_2 \smallstep^* e'
\]
\end{theorem}

\begin{proof}
Immediate from determinism (Theorem~\ref{thm:determinism}).
\end{proof}

\chapter{Denotational Semantics (Expanded)}
\label{ch:denotational-expanded}

\begin{tcolorbox}[colback=green!5!white,colframe=green!75!black,title=\vexp]
This chapter expands the v6.1 denotational semantics with complete semantic domains for all v7.0 constructs and soundness proofs.
\end{tcolorbox}

\section{Semantic Domains}

\begin{definition}[Complete TKS Semantic Domains]
\label{def:semantic-domains}
The semantic domains for TKS v7.0:

\textbf{Base Domains:}
\begin{align}
\sem{\mathsf{Int}} &= \mathbb{Z} \\
\sem{\mathsf{Bool}} &= \mathbb{B} = \{\mathsf{true}, \mathsf{false}\} \\
\sem{\mathsf{Unit}} &= \{()\}
\end{align}

\textbf{TKS-Specific Domains:}
\begin{align}
\sem{\mathsf{World}} &= \{A, B, C, D\} \\
\sem{\mathsf{Element}} &= \sem{\mathsf{World}} \times \{1, \ldots, 10\} \\
\sem{\mathsf{Domain}} &= \{D_1, \ldots, D_7\} \\
\sem{\mathsf{Foundation}} &= \sem{\mathsf{Domain}} \times \{a, b, c, d\} \\
\sem{\mathsf{Noetic}} &= \{\noe{0}, \noe{1}, \ldots, \noe{9}\}
\end{align}

\textbf{Fractal Domain:}
\[
\sem{\mathsf{Frac}[\bar{k}]} = \sem{\mathsf{Domain}} \to \sem{\mathsf{Domain}}
\]

\textbf{RPM Domain:}
\[
\sem{\mathsf{RPM}[\tau]} = \mathcal{P}(\Acquisitions) \to (\sem{\tau} + \mathsf{Failure}) \times \mathcal{P}(\Acquisitions)
\]

\textbf{Function Domain:}
\[
\sem{\tau_1 \to \tau_2} = \sem{\tau_1} \to \sem{\tau_2}
\]

\textbf{Product and Sum:}
\begin{align}
\sem{\tau_1 \times \tau_2} &= \sem{\tau_1} \times \sem{\tau_2} \\
\sem{\tau_1 + \tau_2} &= \sem{\tau_1} + \sem{\tau_2}
\end{align}
\end{definition}

\section{Extended Semantic Function}

\begin{definition}[Denotational Semantics $\sem{\cdot}$]
\label{def:denotation-extended}

\textbf{Variables:}
\[
\sem{x}_\rho = \rho(x)
\]

\textbf{Abstractions:}
\[
\sem{\lambda x:\tau. e}_\rho = \lambda v \in \sem{\tau}. \sem{e}_{\rho[x \mapsto v]}
\]

\textbf{Applications:}
\[
\sem{e_1\; e_2}_\rho = \sem{e_1}_\rho(\sem{e_2}_\rho)
\]

\textbf{Noetic Operations:}
\[
\sem{\noe{k}(e)}_\rho = \mathfrak{n}_k(\sem{e}_\rho)
\]
where $\mathfrak{n}_k : \sem{\mathsf{Domain}} \to \sem{\mathsf{Domain}}$ is the semantic noetic function.

\textbf{Fractal Application:}
\[
\sem{\langle k_1 : \cdots : k_n \rangle(e)}_\rho = (\mathfrak{n}_{k_1} \circ \cdots \circ \mathfrak{n}_{k_n})(\sem{e}_\rho)
\]

\textbf{RPM Return:}
\[
\sem{\rpmret\; e}_\rho = \lambda \alpha. (\mathsf{Success}(\sem{e}_\rho), \alpha)
\]

\textbf{RPM Bind:}
\[
\sem{m \rpmbind f}_\rho = \lambda \alpha.
\begin{cases}
\sem{f}_\rho(v)(\alpha') & \text{if } \sem{m}_\rho(\alpha) = (\mathsf{Success}(v), \alpha') \\
(\mathsf{Failure}(p), \alpha') & \text{if } \sem{m}_\rho(\alpha) = (\mathsf{Failure}(p), \alpha')
\end{cases}
\]

\textbf{Prerequisite Check:}
\[
\sem{\mathsf{check}(p)}_\rho = \lambda \alpha.
\begin{cases}
(\mathsf{Success}(()), \alpha) & \text{if } p \in \alpha \\
(\mathsf{Failure}(p), \alpha) & \text{otherwise}
\end{cases}
\]

\textbf{Acquisition:}
\[
\sem{\mathsf{acquire}(p)}_\rho = \lambda \alpha. (\mathsf{Success}(()), \alpha \cup \{p\})
\]

\textbf{ACBE:}
\[
\sem{\mathrm{ACBE}(e)}_\rho = \mathsf{cascade}(\sem{e}_\rho)
\]
where $\mathsf{cascade}(An) = Dn$.
\end{definition}

\section{Soundness Theorem}

\begin{theorem}[Semantic Soundness]
\label{thm:semantic-soundness}
The operational and denotational semantics agree:
\[
\langle e, \mathcal{E} \rangle \bigstep (v, \mathcal{E}') \implies \sem{e}_{\rho_\mathcal{E}} = \sem{v}_{\rho_{\mathcal{E}'}}
\]
where $\rho_\mathcal{E}$ extracts the variable environment from $\mathcal{E}$.
\end{theorem}

\begin{proof}
By structural induction on the big-step derivation.

\textbf{Base Cases:}

\textit{Case [E-Element]:} $\langle Wn, \mathcal{E} \rangle \bigstep (Wn, \mathcal{E})$
\[
\sem{Wn}_\rho = (W, n) = \sem{Wn}_\rho \quad \checkmark
\]

\textit{Case [E-Foundation]:} Similar.

\textbf{Inductive Cases:}

\textit{Case [E-Noe]:}
\[
\frac{\langle e, \mathcal{E} \rangle \bigstep (v, \mathcal{E}')}
{\langle \noe{k}(e), \mathcal{E} \rangle \bigstep (\mathfrak{n}_k(v), \mathcal{E}')}
\]
By IH: $\sem{e}_\rho = \sem{v}_{\rho'}$

Then:
\[
\sem{\noe{k}(e)}_\rho = \mathfrak{n}_k(\sem{e}_\rho) = \mathfrak{n}_k(\sem{v}_{\rho'}) = \sem{\mathfrak{n}_k(v)}_{\rho'} \quad \checkmark
\]

\textit{Case [E-Frac-N]:} By induction on fractal length, using the IH and compositional structure of fractal denotation.

\textit{Case [S-Ret-Bind]:}
\[
(\rpmret\; v) \rpmbind f \smallstep f(v)
\]
\begin{align}
\sem{(\rpmret\; v) \rpmbind f}_\rho(\alpha) &= \sem{f}_\rho(v)(\alpha) \\
&= \sem{f(v)}_\rho(\alpha) \quad \checkmark
\end{align}

All other cases follow similarly by structural induction.
\end{proof}

\begin{theorem}[Adequacy]
\label{thm:adequacy}
Denotational equality implies observational equivalence:
\[
\sem{e_1}_\rho = \sem{e_2}_\rho \implies e_1 \approx_{obs} e_2
\]
\end{theorem}

\begin{proof}
By showing that denotationally equal programs produce the same observable results in all contexts, using the compositionality of $\sem{\cdot}$.
\end{proof}

%% ============================================================================
%% PART VI: TYPE THEORY (EXPANDED)
%% ============================================================================
\part{Complete Type Theory}

\chapter{Complete Type System}
\label{ch:complete-types}

\begin{tcolorbox}[colback=green!5!white,colframe=green!75!black,title=\vexp]
Complete type system for TKS v7.0 with all type constructors, inference rules, polymorphism, and subtyping.
\end{tcolorbox}

\section{Type Syntax}

\begin{definition}[TKS Type Grammar]
\label{def:type-grammar}
\begin{align}
\tau ::=&\; \mathsf{Int} \mid \mathsf{Bool} \mid \mathsf{Unit} \quad \text{(base types)} \\
    &\mid \mathsf{Domain} \mid \mathsf{Foundation} \mid \mathsf{Aspect} \quad \text{(TKS types)} \\
    &\mid \mathsf{Element}[W] \quad \text{(world-indexed element)} \\
    &\mid \mathsf{Noe}[k] \quad \text{(noetic function type, } k \in \{0,\ldots,9\}) \\
    &\mid \mathsf{Frac}[\langle k_1:\cdots:k_n \rangle] \quad \text{(fractal type)} \\
    &\mid \mathsf{RPM}[\tau] \quad \text{(RPM monad type)} \\
    &\mid \tau_1 \to \tau_2 \quad \text{(function type)} \\
    &\mid \tau_1 \times \tau_2 \quad \text{(product type)} \\
    &\mid \tau_1 + \tau_2 \quad \text{(sum type)} \\
    &\mid \forall \alpha. \tau \quad \text{(universal type)} \\
    &\mid \alpha \quad \text{(type variable)}
\end{align}
\end{definition}

\section{Subtyping}

\begin{definition}[Subtyping Relation $\subtype$]
\label{def:subtyping}

\textbf{Reflexivity:}
\[
\frac{}{\tau \subtype \tau} \quad [\textsc{Sub-Refl}]
\]

\textbf{Transitivity:}
\[
\frac{\tau_1 \subtype \tau_2 \quad \tau_2 \subtype \tau_3}{\tau_1 \subtype \tau_3} \quad [\textsc{Sub-Trans}]
\]

\textbf{Domain Hierarchy:}
\[
\frac{}{\mathsf{Aspect} \subtype \mathsf{Foundation}} \quad [\textsc{Sub-Asp-Fnd}]
\]
\[
\frac{}{\mathsf{Foundation} \subtype \mathsf{Domain}} \quad [\textsc{Sub-Fnd-Dom}]
\]

\textbf{Function Contravariance/Covariance:}
\[
\frac{\tau_1' \subtype \tau_1 \quad \tau_2 \subtype \tau_2'}{\tau_1 \to \tau_2 \subtype \tau_1' \to \tau_2'} \quad [\textsc{Sub-Arrow}]
\]

\textbf{RPM Covariance:}
\[
\frac{\tau \subtype \tau'}{\mathsf{RPM}[\tau] \subtype \mathsf{RPM}[\tau']} \quad [\textsc{Sub-RPM}]
\]

\textbf{Product Covariance:}
\[
\frac{\tau_1 \subtype \tau_1' \quad \tau_2 \subtype \tau_2'}{\tau_1 \times \tau_2 \subtype \tau_1' \times \tau_2'} \quad [\textsc{Sub-Prod}]
\]
\end{definition}

\begin{definition}[Subtyping Lattice]
\label{def:subtype-lattice}
The TKS type hierarchy forms a lattice:
\begin{center}
\begin{tikzpicture}[node distance=1.5cm]
  \node (any) {$\top$ (Any)};
  \node (dom) [below left of=any] {$\mathsf{Domain}$};
  \node (rpm) [below of=any] {$\mathsf{RPM}[\tau]$};
  \node (frac) [below right of=any] {$\mathsf{Frac}[\bar{k}]$};
  \node (fnd) [below of=dom] {$\mathsf{Foundation}$};
  \node (asp) [below of=fnd] {$\mathsf{Aspect}$};
  \node (bot) [below of=asp] {$\bot$ (Nothing)};

  \draw[->] (asp) -- (fnd);
  \draw[->] (fnd) -- (dom);
  \draw[->] (dom) -- (any);
  \draw[->] (rpm) -- (any);
  \draw[->] (frac) -- (any);
  \draw[->] (bot) -- (asp);
\end{tikzpicture}
\end{center}
\end{definition}

\section{Complete Typing Rules}

\begin{definition}[Typing Judgment $\Gamma \vdash e : \tau$]
\label{def:typing-rules-complete}

\textbf{Variables:}
\[
\frac{x : \tau \in \Gamma}{\Gamma \vdash x : \tau} \quad [\textsc{T-Var}]
\]

\textbf{Abstraction:}
\[
\frac{\Gamma, x : \tau_1 \vdash e : \tau_2}{\Gamma \vdash \lambda x:\tau_1. e : \tau_1 \to \tau_2} \quad [\textsc{T-Abs}]
\]

\textbf{Application:}
\[
\frac{\Gamma \vdash e_1 : \tau_1 \to \tau_2 \quad \Gamma \vdash e_2 : \tau_1}{\Gamma \vdash e_1\; e_2 : \tau_2} \quad [\textsc{T-App}]
\]

\textbf{Subsumption:}
\[
\frac{\Gamma \vdash e : \tau \quad \tau \subtype \tau'}{\Gamma \vdash e : \tau'} \quad [\textsc{T-Sub}]
\]

\textbf{Elements:}
\[
\frac{W \in \{A, B, C, D\} \quad n \in \{1, \ldots, 10\}}{\Gamma \vdash Wn : \mathsf{Element}[W]} \quad [\textsc{T-Elem}]
\]

\textbf{Foundations:}
\[
\frac{m \in \{1, \ldots, 7\} \quad w \in \{a, b, c, d\}}{\Gamma \vdash F_{mw} : \mathsf{Foundation}} \quad [\textsc{T-Fnd}]
\]

\textbf{Noetic Application:}
\[
\frac{\Gamma \vdash e : \mathsf{Domain} \quad k \in \{0, \ldots, 9\}}{\Gamma \vdash \noe{k}(e) : \mathsf{Domain}} \quad [\textsc{T-Noe}]
\]

\textbf{Fractal Application:}
\[
\frac{\Gamma \vdash e : \mathsf{Domain} \quad \forall i. k_i \in \{0, \ldots, 9\}}{\Gamma \vdash \langle k_1 : \cdots : k_n \rangle(e) : \mathsf{Domain}} \quad [\textsc{T-Frac}]
\]

\textbf{RPM Return:}
\[
\frac{\Gamma \vdash e : \tau}{\Gamma \vdash \rpmret\; e : \mathsf{RPM}[\tau]} \quad [\textsc{T-Ret}]
\]

\textbf{RPM Bind:}
\[
\frac{\Gamma \vdash m : \mathsf{RPM}[\tau_1] \quad \Gamma \vdash f : \tau_1 \to \mathsf{RPM}[\tau_2]}{\Gamma \vdash m \rpmbind f : \mathsf{RPM}[\tau_2]} \quad [\textsc{T-Bind}]
\]

\textbf{Check:}
\[
\frac{p : \mathsf{Prereq}}{\Gamma \vdash \mathsf{check}(p) : \mathsf{RPM}[\mathsf{Unit}]} \quad [\textsc{T-Check}]
\]

\textbf{Acquire:}
\[
\frac{p : \mathsf{Prereq}}{\Gamma \vdash \mathsf{acquire}(p) : \mathsf{RPM}[\mathsf{Unit}]} \quad [\textsc{T-Acq}]
\]

\textbf{Type Abstraction:}
\[
\frac{\Gamma, \alpha \vdash e : \tau}{\Gamma \vdash \Lambda \alpha. e : \forall \alpha. \tau} \quad [\textsc{T-TAbs}]
\]

\textbf{Type Application:}
\[
\frac{\Gamma \vdash e : \forall \alpha. \tau}{\Gamma \vdash e[\tau'] : \tau[\tau'/\alpha]} \quad [\textsc{T-TApp}]
\]
\end{definition}

\section{Type Inference}

\begin{definition}[Constraint-Based Type Inference]
\label{def:type-inference}
Algorithm $\mathcal{W}$ generates constraints and solves via unification:

\textbf{Constraint Generation:}
\begin{align}
\mathcal{C}(x, \Gamma) &= (\Gamma(x), \emptyset) \\
\mathcal{C}(\lambda x. e, \Gamma) &= \text{let } \alpha \text{ fresh}, (\tau, C) = \mathcal{C}(e, \Gamma[x \mapsto \alpha]) \\
&\quad \text{in } (\alpha \to \tau, C) \\
\mathcal{C}(e_1\; e_2, \Gamma) &= \text{let } (\tau_1, C_1) = \mathcal{C}(e_1, \Gamma), \\
&\quad (\tau_2, C_2) = \mathcal{C}(e_2, \Gamma), \alpha \text{ fresh} \\
&\quad \text{in } (\alpha, C_1 \cup C_2 \cup \{\tau_1 = \tau_2 \to \alpha\})
\end{align}

\textbf{Unification:}
$\mathsf{unify}(C)$ produces most general substitution $\theta$ such that $\theta(C)$ holds.
\end{definition}

\chapter{Progress and Preservation Proofs}
\label{ch:type-proofs}

\begin{tcolorbox}[colback=blue!5!white,colframe=blue!75!black,title=\vnew]
Complete proofs of type safety: Progress and Preservation theorems.
\end{tcolorbox}

\section{Canonical Forms}

\begin{lemma}[Canonical Forms]
\label{lem:canonical}
If $\vdash v : \tau$ and $v$ is a value, then:
\begin{enumerate}
  \item If $\tau = \mathsf{Bool}$, then $v = \mathsf{true}$ or $v = \mathsf{false}$
  \item If $\tau = \mathsf{Int}$, then $v = n$ for some $n \in \mathbb{Z}$
  \item If $\tau = \tau_1 \to \tau_2$, then $v = \lambda x:\tau_1. e$ for some $x, e$
  \item If $\tau = \mathsf{Domain}$, then $v = Wn$ or $v = F_{mw}$ for appropriate indices
  \item If $\tau = \mathsf{RPM}[\tau']$, then $v = \rpmret\; v'$ or $v = \Failed(p)$
  \item If $\tau = \tau_1 \times \tau_2$, then $v = (v_1, v_2)$ for some $v_1, v_2$
\end{enumerate}
\end{lemma}

\begin{proof}
By inspection of the typing rules and value forms. Each type has a unique set of value constructors, and well-typed values must use the appropriate constructor.
\end{proof}

\section{Progress Theorem}

\begin{theorem}[Progress]
\label{thm:progress}
If $\vdash e : \tau$, then either:
\begin{enumerate}
  \item $e$ is a value, or
  \item There exists $e'$ such that $e \smallstep e'$
\end{enumerate}
\end{theorem}

\begin{proof}
By structural induction on the typing derivation $\vdash e : \tau$.

\textbf{Case [T-Var]:} Cannot occur in empty context.

\textbf{Case [T-Abs]:} $e = \lambda x:\tau_1. e'$ is already a value.

\textbf{Case [T-App]:} $e = e_1\; e_2$ where $\vdash e_1 : \tau_1 \to \tau_2$ and $\vdash e_2 : \tau_1$.

By IH on $e_1$:
\begin{itemize}
  \item If $e_1 \smallstep e_1'$, then $e_1\; e_2 \smallstep e_1'\; e_2$ by [S-App-L].
  \item If $e_1$ is a value, by IH on $e_2$:
    \begin{itemize}
      \item If $e_2 \smallstep e_2'$, then $e_1\; e_2 \smallstep e_1\; e_2'$ by [S-App-R].
      \item If $e_2$ is a value, by Canonical Forms, $e_1 = \lambda x:\tau_1. e'$.
        Then $(\lambda x:\tau_1. e')\; e_2 \smallstep e'[e_2/x]$ by [S-Beta].
    \end{itemize}
\end{itemize}

\textbf{Case [T-Noe]:} $e = \noe{k}(e')$ where $\vdash e' : \mathsf{Domain}$.

By IH on $e'$:
\begin{itemize}
  \item If $e' \smallstep e''$, then $\noe{k}(e') \smallstep \noe{k}(e'')$ by [S-Noe-Cong].
  \item If $e'$ is a value $v$, then $\noe{k}(v) \smallstep \mathfrak{n}_k(v)$ by [S-Noe].
\end{itemize}

\textbf{Case [T-Frac]:} Similar to [T-Noe], using fractal reduction rules.

\textbf{Case [T-Ret]:} $e = \rpmret\; e'$.
\begin{itemize}
  \item If $e' \smallstep e''$, then $\rpmret\; e' \smallstep \rpmret\; e''$.
  \item If $e'$ is a value, then $\rpmret\; e'$ is a value.
\end{itemize}

\textbf{Case [T-Bind]:} $e = m \rpmbind f$ where $\vdash m : \mathsf{RPM}[\tau_1]$.

By IH on $m$:
\begin{itemize}
  \item If $m \smallstep m'$, then $m \rpmbind f \smallstep m' \rpmbind f$ by [S-Bind-Cong].
  \item If $m$ is a value, by Canonical Forms:
    \begin{itemize}
      \item If $m = \rpmret\; v$, then $(\rpmret\; v) \rpmbind f \smallstep f(v)$ by [S-Ret-Bind].
      \item If $m = \Failed(p)$, then $\Failed(p) \rpmbind f \smallstep \Failed(p)$ by [S-Fail-Bind].
    \end{itemize}
\end{itemize}

\textbf{Case [T-Check]:} $e = \mathsf{check}(p)$.
This steps to either $\rpmret\; ()$ or $\Failed(p)$ depending on state.

\textbf{Case [T-Acq]:} $e = \mathsf{acquire}(p)$.
This steps to $\rpmret\; ()$ with updated state.

All other cases follow similarly.
\end{proof}

\section{Preservation Theorem}

\begin{lemma}[Substitution]
\label{lem:substitution}
If $\Gamma, x : \tau' \vdash e : \tau$ and $\Gamma \vdash v : \tau'$, then $\Gamma \vdash e[v/x] : \tau$.
\end{lemma}

\begin{proof}
By structural induction on the derivation of $\Gamma, x : \tau' \vdash e : \tau$.
\end{proof}

\begin{theorem}[Preservation]
\label{thm:preservation}
If $\Gamma \vdash e : \tau$ and $e \smallstep e'$, then $\Gamma \vdash e' : \tau$.
\end{theorem}

\begin{proof}
By structural induction on the derivation of $e \smallstep e'$.

\textbf{Case [S-Beta]:} $(\lambda x:\tau_1. e)\; v \smallstep e[v/x]$

From premises: $\Gamma \vdash (\lambda x:\tau_1. e)\; v : \tau_2$

By inversion on [T-App]:
\begin{itemize}
  \item $\Gamma \vdash \lambda x:\tau_1. e : \tau_1 \to \tau_2$
  \item $\Gamma \vdash v : \tau_1$
\end{itemize}

By inversion on [T-Abs]: $\Gamma, x : \tau_1 \vdash e : \tau_2$

By Substitution Lemma: $\Gamma \vdash e[v/x] : \tau_2$ \checkmark

\textbf{Case [S-Noe]:} $\noe{k}(v) \smallstep \mathfrak{n}_k(v)$

From premises: $\Gamma \vdash \noe{k}(v) : \mathsf{Domain}$

By [T-Noe]: $\Gamma \vdash v : \mathsf{Domain}$

The semantic function $\mathfrak{n}_k$ preserves the Domain type:
$\Gamma \vdash \mathfrak{n}_k(v) : \mathsf{Domain}$ \checkmark

\textbf{Case [S-Frac-1]:} $\langle k \rangle(v) \smallstep \noe{k}(v)$

From premises: $\Gamma \vdash \langle k \rangle(v) : \mathsf{Domain}$

By [T-Frac]: $\Gamma \vdash v : \mathsf{Domain}$

By [T-Noe]: $\Gamma \vdash \noe{k}(v) : \mathsf{Domain}$ \checkmark

\textbf{Case [S-Ret-Bind]:} $(\rpmret\; v) \rpmbind f \smallstep f(v)$

From premises: $\Gamma \vdash (\rpmret\; v) \rpmbind f : \mathsf{RPM}[\tau_2]$

By inversion on [T-Bind]:
\begin{itemize}
  \item $\Gamma \vdash \rpmret\; v : \mathsf{RPM}[\tau_1]$
  \item $\Gamma \vdash f : \tau_1 \to \mathsf{RPM}[\tau_2]$
\end{itemize}

By inversion on [T-Ret]: $\Gamma \vdash v : \tau_1$

By [T-App]: $\Gamma \vdash f(v) : \mathsf{RPM}[\tau_2]$ \checkmark

\textbf{Case [S-Fail-Bind]:} $\Failed(p) \rpmbind f \smallstep \Failed(p)$

$\Failed(p) : \mathsf{RPM}[\tau]$ for any $\tau$, including the result type. \checkmark

\textbf{Case [S-Check-Ok]:} $\mathsf{check}(p) \smallstep \rpmret\; ()$

From [T-Check]: $\Gamma \vdash \mathsf{check}(p) : \mathsf{RPM}[\mathsf{Unit}]$

By [T-Ret]: $\Gamma \vdash \rpmret\; () : \mathsf{RPM}[\mathsf{Unit}]$ \checkmark

\textbf{Case [S-Acquire]:} $\mathsf{acquire}(p) \smallstep \rpmret\; ()$

Similar to [S-Check-Ok]. \checkmark

\textbf{Congruence Cases:} Follow by IH on the sub-derivation.
\end{proof}

\begin{theorem}[Type Safety]
\label{thm:type-safety}
Well-typed TKS programs do not get stuck:
\[
\text{If } \vdash e : \tau, \text{ then } e \text{ evaluates to a value or diverges.}
\]
\end{theorem}

\begin{proof}
By Progress (Theorem~\ref{thm:progress}) and Preservation (Theorem~\ref{thm:preservation}).

If $\vdash e : \tau$ and $e$ is not a value, Progress guarantees $e \smallstep e'$ for some $e'$. Preservation guarantees $\vdash e' : \tau$. By induction, either evaluation terminates at a value or continues indefinitely.
\end{proof}

%% ============================================================================
%% PART VII: ADVANCED TOPICS (NEW)
%% ============================================================================
\part{Advanced Topics}

\chapter{Concurrency Model (Prototype)}
\label{ch:concurrency}

\begin{tcolorbox}[colback=blue!5!white,colframe=blue!75!black,title=\vnew]
This chapter introduces a prototype concurrency model for TKS, enabling parallel fractal evaluation and synchronized transformations.
\end{tcolorbox}

\section{Motivation}

Real-world transformation processes often involve parallel development across multiple domains or simultaneous work on independent aspects. The TKS concurrency model formalizes:
\begin{itemize}
  \item Parallel evaluation of independent fractals
  \item Synchronization when fractals share state
  \item Determinism conditions for predictable outcomes
  \item Resource management for acquisitions
\end{itemize}

\section{Parallel Fractal Composition}

\begin{definition}[Independence]
\label{def:independence}
Two fractals $\Frac_1$ and $\Frac_2$ are \textbf{independent} if their prerequisite sets are disjoint:
\[
\mathsf{prereqs}(\Frac_1) \cap \mathsf{prereqs}(\Frac_2) = \emptyset
\]
where $\mathsf{prereqs}(\Frac)$ is the set of acquisitions checked or modified by $\Frac$.
\end{definition}

\begin{definition}[Parallel Composition $\pcomp$]
\label{def:parallel-comp}
For independent fractals $\Frac_1 : \mathsf{Domain} \to \mathsf{Domain}$ and $\Frac_2 : \mathsf{Domain} \to \mathsf{Domain}$:
\[
\Frac_1 \pcomp \Frac_2 : \mathsf{Domain} \times \mathsf{Domain} \to \mathsf{Domain} \times \mathsf{Domain}
\]
\[
(\Frac_1 \pcomp \Frac_2)(E_1, E_2) = (\Frac_1(E_1), \Frac_2(E_2))
\]
\end{definition}

\begin{definition}[Parallel RPM Execution]
\label{def:parallel-rpm}
For RPM computations $m_1 : \mathsf{RPM}[\tau_1]$ and $m_2 : \mathsf{RPM}[\tau_2]$:
\[
m_1 \pcomp m_2 : \mathsf{RPM}[\tau_1 \times \tau_2]
\]

\textbf{Semantics:}
\[
(m_1 \pcomp m_2)(\alpha) =
\begin{cases}
(\mathsf{Success}(v_1, v_2), \alpha_1 \cup \alpha_2) & \text{if both succeed} \\
(\mathsf{Failure}(p), \alpha') & \text{if either fails first}
\end{cases}
\]
where $m_1(\alpha) = (\mathsf{Success}(v_1), \alpha_1)$ and $m_2(\alpha) = (\mathsf{Success}(v_2), \alpha_2)$.
\end{definition}

\section{Synchronization Primitives}

\begin{definition}[Fractal Barrier]
\label{def:fractal-barrier}
A \textbf{barrier} synchronizes parallel fractal threads:
\[
\mathsf{barrier} : \mathsf{RPM}[\tau_1] \times \mathsf{RPM}[\tau_2] \to \mathsf{RPM}[\tau_1 \times \tau_2]
\]

The barrier waits for both computations to complete (or fail) before proceeding:
\[
\mathsf{barrier}(m_1, m_2)(\alpha) = (m_1 \pcomp m_2)(\alpha)
\]
\end{definition}

\begin{definition}[Acquisition Lock]
\label{def:acq-lock}
An \textbf{acquisition lock} prevents concurrent modification:
\[
\mathsf{lock}(p) : \mathsf{RPM}[\mathsf{Unit}]
\]
\[
\mathsf{unlock}(p) : \mathsf{RPM}[\mathsf{Unit}]
\]

\textbf{Semantics:}
\[
\mathsf{lock}(p)(\alpha) =
\begin{cases}
(\mathsf{Success}(()), \alpha \cup \{\mathsf{locked}_p\}) & \text{if } \mathsf{locked}_p \notin \alpha \\
(\mathsf{Blocked}(p), \alpha) & \text{otherwise}
\end{cases}
\]
\end{definition}

\begin{definition}[Fractal Channel]
\label{def:fractal-channel}
A \textbf{channel} enables communication between parallel fractal evaluations:
\[
\mathsf{Channel}[\tau] = (\mathsf{send} : \tau \to \mathsf{RPM}[\mathsf{Unit}], \mathsf{recv} : \mathsf{RPM}[\tau])
\]
\end{definition}

\section{Determinacy Analysis}

\begin{theorem}[Confluence for Independent Fractals]
\label{thm:parallel-confluence}
If $\Frac_1$ and $\Frac_2$ are independent, then:
\[
(\Frac_1 \circ \Frac_2)(E) = (\Frac_2 \circ \Frac_1)(E)
\]
Parallel evaluation produces the same result regardless of interleaving.
\end{theorem}

\begin{proof}
Independence ensures no shared state modifications. Each fractal operates on disjoint portions of the acquisition state, so the final state is the union of individual effects, which is commutative.
\end{proof}

\begin{definition}[Conflict Detection]
\label{def:conflict-detection}
A \textbf{conflict} occurs when two concurrent operations access the same acquisition with at least one write:
\[
\mathsf{conflict}(m_1, m_2) \iff \mathsf{writes}(m_1) \cap \mathsf{accesses}(m_2) \neq \emptyset
\]
\end{definition}

\begin{theorem}[Race-Free Implies Deterministic]
\label{thm:race-free}
If a parallel TKS program has no conflicts, its evaluation is deterministic.
\end{theorem}

\section{Parallel Evaluation Rules}

\begin{definition}[Parallel Small-Step]
\label{def:parallel-smallstep}

\textbf{Parallel Step (Both):}
\[
\frac{e_1 \smallstep e_1' \quad e_2 \smallstep e_2'}
{e_1 \pcomp e_2 \smallstep e_1' \pcomp e_2'} \quad [\textsc{P-Both}]
\]

\textbf{Parallel Step (Left):}
\[
\frac{e_1 \smallstep e_1'}
{e_1 \pcomp e_2 \smallstep e_1' \pcomp e_2} \quad [\textsc{P-Left}]
\]

\textbf{Parallel Step (Right):}
\[
\frac{e_2 \smallstep e_2'}
{e_1 \pcomp e_2 \smallstep e_1 \pcomp e_2'} \quad [\textsc{P-Right}]
\]

\textbf{Parallel Join:}
\[
\frac{}
{v_1 \pcomp v_2 \smallstep (v_1, v_2)} \quad [\textsc{P-Join}]
\]
\end{definition}

\section{Example: Parallel Domain Transformation}

\begin{example}[Parallel Health and Wealth Transformation]
\label{ex:parallel-transform}
Transform Health (D2) and Wealth (D6) domains simultaneously:

\begin{verbatim}
let parallel_transform : RPM[Domain * Domain] =
  let health_work =
    check(D2) >>= \_ . return (<1:4:7>(D2_health))
  in
  let wealth_work =
    check(D6) >>= \_ . return (<1:4:7>(D6_wealth))
  in
  health_work || wealth_work
\end{verbatim}

\textbf{Execution:} Both transformations proceed concurrently since D2 and D6 are independent domains.

\textbf{Result:} $((\text{MVR}(D2_{\text{health}}), \text{MVR}(D6_{\text{wealth}})), \alpha)$
\end{example}

\chapter{Quantum Noetic Algebra (Prototype)}
\label{ch:quantum}

\begin{tcolorbox}[colback=blue!5!white,colframe=blue!75!black,title=\vnew]
This chapter introduces a quantum-mechanical interpretation of TKS, treating Noetics as operators on a Hilbert space of Ideas.
\end{tcolorbox}

\section{Motivation}

The quantum interpretation of TKS provides:
\begin{itemize}
  \item Superposition of Ideas (mixed states)
  \item Probabilistic transformation outcomes
  \item Non-commutative Noetic effects
  \item Entanglement between related Ideas
\end{itemize}

\section{Quantum Idea Space}

\begin{definition}[Hilbert Space of Ideas]
\label{def:hilbert-ideas}
The \textbf{quantum Idea space} is the Hilbert space:
\[
\Hspace_\Ideas = \ell^2(\Ideas) = \left\{ \psi : \Ideas \to \mathbb{C} \mid \sum_{E \in \Ideas} |\psi(E)|^2 < \infty \right\}
\]
with inner product:
\[
\langle \psi | \phi \rangle = \sum_{E \in \Ideas} \overline{\psi(E)} \phi(E)
\]
\end{definition}

\begin{definition}[Idea Basis States]
\label{def:idea-basis}
For each classical Idea $E \in \Ideas$, there is a basis state:
\[
\qket{E} \in \Hspace_\Ideas
\]
The basis states are orthonormal:
\[
\qbraket{E}{E'} = \delta_{E,E'}
\]
\end{definition}

\begin{definition}[Quantum Idea State]
\label{def:quantum-state}
A \textbf{quantum Idea state} is a normalized vector:
\[
\qket{\psi} = \sum_{E \in \Ideas} c_E \qket{E}, \quad \sum_E |c_E|^2 = 1
\]
where $c_E \in \mathbb{C}$ are probability amplitudes.
\end{definition}

\section{Noetics as Quantum Operators}

\begin{definition}[Noetic Operator]
\label{def:noetic-operator}
Each Noetic $\noe{k}$ becomes a linear operator:
\[
\qop{\nu}_k : \Hspace_\Ideas \to \Hspace_\Ideas
\]
defined by its action on basis states:
\[
\qop{\nu}_k \qket{E} = \qket{\mathfrak{n}_k(E)}
\]
and extended linearly to superpositions.
\end{definition}

\begin{proposition}[Noetic Operators are Unitary]
\label{prop:noetic-unitary}
If $\mathfrak{n}_k$ is a bijection on $\Ideas$, then $\qop{\nu}_k$ is unitary:
\[
\qop{\nu}_k^\dagger \qop{\nu}_k = \qop{\nu}_k \qop{\nu}_k^\dagger = \mathbf{I}
\]
\end{proposition}

\begin{proof}
For bijective $\mathfrak{n}_k$:
\[
\langle \qop{\nu}_k \psi | \qop{\nu}_k \phi \rangle = \sum_E \overline{(\qop{\nu}_k \psi)(E)} (\qop{\nu}_k \phi)(E) = \sum_E \overline{\psi(\mathfrak{n}_k^{-1}(E))} \phi(\mathfrak{n}_k^{-1}(E)) = \langle \psi | \phi \rangle
\]
\end{proof}

\section{Superposition and Measurement}

\begin{definition}[Superposition of Ideas]
\label{def:superposition}
A superposition represents uncertainty about which Idea is present:
\[
\qket{\psi} = \frac{1}{\sqrt{2}}(\qket{A1} + \qket{B1})
\]
This state is ``half Spiritual, half Mental'' for Idea 1.
\end{definition}

\begin{definition}[Measurement Postulate]
\label{def:measurement}
Measuring Noetic $\qop{\nu}_k$ on state $\qket{\psi}$ yields eigenvalue $\lambda$ with probability:
\[
P(\lambda) = |\langle \lambda | \psi \rangle|^2
\]
where $\qket{\lambda}$ is the eigenstate for eigenvalue $\lambda$.

After measurement, the state collapses:
\[
\qket{\psi} \xrightarrow{\text{measure}} \qket{\lambda}
\]
\end{definition}

\begin{example}[World Measurement]
\label{ex:world-measurement}
Given superposition $\qket{\psi} = \frac{1}{2}(\qket{A1} + \qket{B1} + \qket{C1} + \qket{D1})$:

Measuring ``which World?'' yields:
\begin{itemize}
  \item $P(A) = P(B) = P(C) = P(D) = \frac{1}{4}$
\end{itemize}

After observing World $A$, state becomes $\qket{A1}$.
\end{example}

\section{Commutation Relations}

\begin{definition}[Noetic Commutator]
\label{def:noetic-commutator}
The commutator of two Noetic operators:
\[
[\qop{\nu}_i, \qop{\nu}_j] = \qop{\nu}_i \qop{\nu}_j - \qop{\nu}_j \qop{\nu}_i
\]
\end{definition}

\begin{theorem}[Noetic Uncertainty Principle]
\label{thm:noetic-uncertainty}
For non-commuting Noetics $[\qop{\nu}_i, \qop{\nu}_j] \neq 0$:
\[
\Delta \nu_i \cdot \Delta \nu_j \geq \frac{1}{2} |\langle [\qop{\nu}_i, \qop{\nu}_j] \rangle|
\]
where $\Delta \nu_k = \sqrt{\langle \qop{\nu}_k^2 \rangle - \langle \qop{\nu}_k \rangle^2}$ is the standard deviation.
\end{theorem}

\begin{corollary}[Above/Below Uncertainty]
\label{cor:above-below}
The Noetics $\noe{2}$ (Above) and $\noe{3}$ (Below) may exhibit complementarity:
\[
\Delta \nu_2 \cdot \Delta \nu_3 \geq \frac{1}{2} |\langle [\qop{\nu}_2, \qop{\nu}_3] \rangle|
\]
Perfect knowledge of ``aboveness'' limits knowledge of ``belowness.''
\end{corollary}

\section{Quantum Fractals}

\begin{definition}[Quantum Fractal]
\label{def:quantum-fractal}
A \textbf{quantum fractal} is the composition of Noetic operators:
\[
\qop{\Frac} = \qop{\nu}_{k_1} \qop{\nu}_{k_2} \cdots \qop{\nu}_{k_n}
\]
For unitary Noetics, $\qop{\Frac}$ is also unitary.
\end{definition}

\begin{definition}[Quantum MVR]
\label{def:quantum-mvr}
The quantum MVR fractal:
\[
\qop{\Frac}_{\text{MVR}} = \qop{\nu}_1 \qop{\nu}_4 \qop{\nu}_7
\]
This is the quantum analogue of conscious-vibration-rhythm transformation.
\end{definition}

\begin{theorem}[Quantum Fractal Evolution]
\label{thm:quantum-evolution}
The time evolution of a quantum Idea state under fractal $\qop{\Frac}$ is:
\[
\qket{\psi(t)} = e^{-i \qop{\Frac} t / \hbar} \qket{\psi(0)}
\]
\end{theorem}

\section{Entanglement}

\begin{definition}[Entangled Ideas]
\label{def:entanglement}
Two Ideas $E_1, E_2$ are \textbf{entangled} if their joint state cannot be factored:
\[
\qket{\Psi_{12}} \neq \qket{\psi_1} \otimes \qket{\psi_2}
\]
\end{definition}

\begin{example}[Bell State of Ideas]
\label{ex:bell-state}
The maximally entangled state:
\[
\qket{\Phi^+} = \frac{1}{\sqrt{2}}(\qket{A1} \otimes \qket{A1} + \qket{D1} \otimes \qket{D1})
\]
Measuring one Idea immediately determines the other---Spiritual-Spiritual or Physical-Physical.
\end{example}

\section{Quantum RPM}

\begin{definition}[Quantum Prerequisite Check]
\label{def:quantum-check}
A quantum prerequisite check measures whether an acquisition is present:
\[
\qop{P}_p = \sum_{E : p \in \mathsf{prereqs}(E)} \qket{E}\qbra{E}
\]
This is a projection operator.
\end{definition}

\begin{theorem}[Quantum Zeno Effect for Prerequisites]
\label{thm:quantum-zeno}
Frequent prerequisite checking can ``freeze'' a quantum state, preventing evolution outside the checked subspace.
\end{theorem}

\section{Interpretation}

The quantum TKS framework suggests:
\begin{enumerate}
  \item \textbf{Ideas exist in superposition} until observed/manifested
  \item \textbf{Noetic operations} are quantum gates transforming Idea states
  \item \textbf{Measurement} (manifestation) collapses superpositions
  \item \textbf{Entanglement} links Ideas across Worlds/Foundations
  \item \textbf{Uncertainty} limits simultaneous knowledge of complementary Noetics
\end{enumerate}

This provides a mathematical framework for understanding phenomena like:
\begin{itemize}
  \item Why focused intention (measurement) affects outcomes
  \item How Ideas in different Worlds can be instantaneously correlated
  \item The role of the observer in manifestation
\end{itemize}

%% ============================================================================
%% PART VIII: LANGUAGE AND IMPLEMENTATION
%% ============================================================================
\part{Language Specification and Implementation}

\chapter{Language Specification}
\label{ch:language-spec}

\section{Abstract Syntax}

\begin{definition}[TKS Expression Grammar]
\label{def:tks-grammar}
\begin{align}
e ::=&\; v \mid x \mid e_1 + e_2 \mid e_1 - e_2 \mid e_1 * e_2 \mid e_1 / e_2 \\
    &\mid \noe{k}(e) \mid \langle k_1 : \cdots : k_n \rangle(e) \\
    &\mid \rpmret\; e \mid e_1 \rpmbind e_2 \\
    &\mid \mathsf{check}(p) \mid \mathsf{acquire}(p) \\
    &\mid \lambda x : \tau. e \mid e_1\; e_2 \\
    &\mid \mathbf{let}\; x = e_1\; \mathbf{in}\; e_2 \\
    &\mid \mathbf{if}\; e_1\; \mathbf{then}\; e_2\; \mathbf{else}\; e_3
\end{align}
\end{definition}

\section{Lexical Conventions}

\begin{itemize}
\item Identifiers: $[a\text{-}zA\text{-}Z][a\text{-}zA\text{-}Z0\text{-}9\_]*$
\item Noetic indices: $0, 1, 2, \ldots, 9$
\item Fractal delimiters: $\langle$ and $\rangle$
\item Fractal separator: $:$
\item Comments: $--$ to end of line
\end{itemize}

\section{Precedence and Associativity}

\begin{center}
\begin{tabular}{cll}
\toprule
\textbf{Level} & \textbf{Operators} & \textbf{Associativity} \\
\midrule
1 (lowest) & $\rpmbind$ & Left \\
2 & $\lambda$, $\mathbf{let}$, $\mathbf{if}$ & Right \\
3 & $+$, $-$ & Left \\
4 & $*$, $/$ & Left \\
5 & Application & Left \\
6 (highest) & $\noe{k}$, $\langle\cdots\rangle$ & Prefix \\
\bottomrule
\end{tabular}
\end{center}

\chapter{Implementation Architecture}
\label{ch:implementation}

\section{Component Overview}

A TKS implementation consists of the following major components:

\begin{enumerate}
\item \textbf{Core Data Structures}: Noetics, Fractals, RPM computations, Domains, Foundations
\item \textbf{Parser}: Converts source text to abstract syntax tree
\item \textbf{Type Checker}: Validates well-typedness and infers types
\item \textbf{Evaluator}: Executes TKS programs according to operational semantics
\item \textbf{RPM Runtime}: Manages acquisition state and prerequisite checking
\item \textbf{Fractal Engine}: Handles fractal composition, simplification, and application
\end{enumerate}

\section{Data Structure Recommendations}

\begin{itemize}
\item \textbf{Noetics}: Enum or ADT with 10 variants
\item \textbf{Fractals}: Non-empty list of noetic indices with simplification on construction
\item \textbf{RPM}: State monad transformer or explicit state-passing
\item \textbf{Acquisition State}: Set or map from acquisitions to boolean
\item \textbf{Types}: Algebraic data type with constructors for base types and type operators
\end{itemize}

\section{Evaluation Strategy}

Recommended: Big-step evaluation with explicit state threading. Small-step reduction available for debugging and step-by-step execution.

\section{Error Handling}

\begin{itemize}
\item Type errors: Reported at compile time with location information
\item RPM failures: Propagated with failure origin for diagnostics
\item Runtime errors: Division by zero, undefined variables (should be caught by type checker)
\end{itemize}

%% ============================================================================
%% PART IX: APPLICATIONS
%% ============================================================================
\part{Applications and Reference}

\chapter{Comprehensive Worked Examples}
\label{ch:examples}

\begin{tcolorbox}[colback=green!5!white,colframe=green!75!black,title=\vexp]
This chapter provides 20+ comprehensive worked examples demonstrating all major TKS constructs with complete evaluation traces.
\end{tcolorbox}

\section{Foundation Examples}

\begin{example}[Complete Seven-Domain Traversal]
\label{ex:seven-domain}
\textbf{Goal:} Transform through all seven domains with prerequisite tracking.

\textbf{Program:}
\begin{verbatim}
let full_traversal : Domain -> RPM[Domain] =
  \initial : Domain .
    acquire(D1) >>= \_ .  -- Spirit
    acquire(D2) >>= \_ .  -- Life
    acquire(D3) >>= \_ .  -- Knowledge
    acquire(D4) >>= \_ .  -- Love
    acquire(D5) >>= \_ .  -- Power
    acquire(D6) >>= \_ .  -- Wealth
    acquire(D7) >>= \_ .  -- Expression
    return (<1:4:7>(initial))
\end{verbatim}

\textbf{Evaluation Trace (initial $\alpha = \emptyset$):}
\begin{align}
&\text{Step 1: } \mathsf{acquire}(D1) \Rightarrow \alpha_1 = \{D1\} \\
&\text{Step 2: } \mathsf{acquire}(D2) \Rightarrow \alpha_2 = \{D1, D2\} \\
&\text{Step 3: } \mathsf{acquire}(D3) \Rightarrow \alpha_3 = \{D1, D2, D3\} \\
&\text{Step 4: } \mathsf{acquire}(D4) \Rightarrow \alpha_4 = \{D1, D2, D3, D4\} \\
&\text{Step 5: } \mathsf{acquire}(D5) \Rightarrow \alpha_5 = \{D1, D2, D3, D4, D5\} \\
&\text{Step 6: } \mathsf{acquire}(D6) \Rightarrow \alpha_6 = \{D1, D2, D3, D4, D5, D6\} \\
&\text{Step 7: } \mathsf{acquire}(D7) \Rightarrow \alpha_7 = \{D1, \ldots, D7\} \\
&\text{Step 8: } \langle 1:4:7 \rangle(\text{initial}) = \noe{1}(\noe{4}(\noe{7}(\text{initial}))) \\
&\text{Final: } (\mathsf{Success}(\text{MVR(initial)}), \{D1, \ldots, D7\})
\end{align}
\end{example}

\begin{example}[Partial Failure Scenario]
\label{ex:partial-failure}
\textbf{Goal:} Demonstrate RPM failure propagation with partial state.

\textbf{Program:}
\begin{verbatim}
let guarded_transform : RPM[Domain] =
  acquire(D1) >>= \_ .
  acquire(D2) >>= \_ .
  check(D5) >>= \_ .    -- Requires Power (not acquired!)
  return (<1:4:7>(D2_health))
\end{verbatim}

\textbf{Evaluation (initial $\alpha = \emptyset$):}
\begin{align}
&\mathsf{acquire}(D1): \alpha_1 = \{D1\} \quad \checkmark \\
&\mathsf{acquire}(D2): \alpha_2 = \{D1, D2\} \quad \checkmark \\
&\mathsf{check}(D5): D5 \notin \alpha_2 \quad \times \\
&\text{Result: } (\mathsf{Failure}(D5), \{D1, D2\})
\end{align}

\textbf{Note:} Acquisitions D1 and D2 persist despite overall failure.
\end{example}

\section{Fractal Application Examples}

\begin{example}[MVR Fractal Step-by-Step]
\label{ex:mvr-stepbystep}
\textbf{Goal:} Trace MVR fractal application in detail.

\textbf{Input:} $E_0 = D3_{2a}$ (Knowledge, Mental, Positive)

\textbf{Small-Step Trace:}
\begin{align}
&\langle 1:4:7 \rangle(D3_{2a}) \\
&\smallstep \noe{1}(\langle 4:7 \rangle(D3_{2a})) \quad [\textsc{S-Frac-Exp}] \\
&\smallstep \noe{1}(\noe{4}(\langle 7 \rangle(D3_{2a}))) \quad [\textsc{S-Frac-Exp}] \\
&\smallstep \noe{1}(\noe{4}(\noe{7}(D3_{2a}))) \quad [\textsc{S-Frac-1}] \\
&\smallstep \noe{1}(\noe{4}(\text{Rhythmic}(D3_{2a}))) \quad [\textsc{S-Noe}] \\
&\smallstep \noe{1}(\text{Vibrating Rhythmic}(D3_{2a})) \quad [\textsc{S-Noe}] \\
&\smallstep \text{Aware Vibrating Rhythmic}(D3_{2a}) \quad [\textsc{S-Noe}]
\end{align}

\textbf{Result:} Conscious, energized, patterned Knowledge.
\end{example}

\begin{example}[Fractal Composition and Simplification]
\label{ex:frac-simplify}
\textbf{Goal:} Demonstrate fractal simplification rules.

\textbf{Setup:}
\begin{align}
\Frac_1 &= \langle 2:3 \rangle \quad \text{(Above-Below)} \\
\Frac_2 &= \langle 3:2 \rangle \quad \text{(Below-Above)}
\end{align}

\textbf{Composition:}
\begin{align}
\Frac_1 \circ \Frac_2 &= \langle 2:3:3:2 \rangle \\
&\Rightarrow \langle 2:0:2 \rangle \quad \text{(dual cancellation: } 3:3 \to 0) \\
&\Rightarrow \langle 2:2 \rangle \quad \text{(identity elimination)} \\
&\Rightarrow \langle 0 \rangle \quad \text{(dual cancellation: } 2:2 \to 0) \\
&= \Frac_{\id}
\end{align}

\textbf{Interpretation:} Above-Below followed by Below-Above returns to identity.
\end{example}

\begin{example}[Eigenfractal Verification]
\label{ex:eigenfrac-verify}
\textbf{Goal:} Find and verify an eigenfractal for MVR.

\textbf{Candidate:} $E^* = \text{Aware Vibrating Rhythmic } A1_{1a}$

\textbf{Verification:}
\begin{align}
\Frac_{\text{MVR}}(E^*) &= \noe{1}(\noe{4}(\noe{7}(E^*))) \\
&= \noe{1}(\noe{4}(\noe{7}(\text{Aware Vibrating Rhythmic } A1_{1a}))) \\
&= \noe{1}(\noe{4}(\text{Aware Vibrating Rhythmic } A1_{1a})) \quad \text{(idempotence of } \noe{7}) \\
&= \noe{1}(\text{Aware Vibrating Rhythmic } A1_{1a}) \quad \text{(idempotence of } \noe{4}) \\
&= \text{Aware Vibrating Rhythmic } A1_{1a} \quad \text{(idempotence of } \noe{1}) \\
&= E^*
\end{align}

\textbf{Conclusion:} $E^*$ is a fixed point of MVR with eigenvalue 1.
\end{example}

\section{RPM Monad Examples}

\begin{example}[Kleisli Composition]
\label{ex:kleisli-comp}
\textbf{Goal:} Demonstrate Kleisli category composition.

\textbf{Functions:}
\begin{align}
f &: \mathsf{Domain} \to \mathsf{RPM}[\mathsf{Domain}] = \lambda d. \mathsf{check}(D1) \rpmbind \lambda\_. \rpmret(\noe{1}(d)) \\
g &: \mathsf{Domain} \to \mathsf{RPM}[\mathsf{Domain}] = \lambda d. \mathsf{check}(D4) \rpmbind \lambda\_. \rpmret(\noe{4}(d))
\end{align}

\textbf{Kleisli Composition:}
\begin{align}
(g \circ_\Kleisli f)(d) &= f(d) \rpmbind g \\
&= (\mathsf{check}(D1) \rpmbind \lambda\_. \rpmret(\noe{1}(d))) \rpmbind g
\end{align}

\textbf{Evaluation with $\alpha = \{D1, D4\}$:}
\begin{align}
&\mathsf{check}(D1)(\alpha) = (\mathsf{Success}(()), \alpha) \\
&\rpmret(\noe{1}(d))(\alpha) = (\mathsf{Success}(\noe{1}(d)), \alpha) \\
&g(\noe{1}(d))(\alpha) = \mathsf{check}(D4) \rpmbind \lambda\_. \rpmret(\noe{4}(\noe{1}(d))) \\
&= (\mathsf{Success}(\noe{4}(\noe{1}(d))), \alpha)
\end{align}

\textbf{Result:} $(\mathsf{Success}(\text{Vibrating Aware } d), \{D1, D4\})$
\end{example}

\begin{example}[RPM Choice Operator]
\label{ex:rpm-choice}
\textbf{Goal:} Demonstrate alternative path selection.

\textbf{Program:}
\begin{verbatim}
let try_paths : RPM[Domain] =
  (check(D5) >>= \_ . return (noe[5](base)))
  <|>
  (check(D6) >>= \_ . return (noe[6](base)))
\end{verbatim}

\textbf{Evaluation with $\alpha = \{D6\}$:}
\begin{enumerate}
  \item Try first alternative: $\mathsf{check}(D5)$
  \item $D5 \notin \alpha$, so first fails
  \item Try second alternative: $\mathsf{check}(D6)$
  \item $D6 \in \alpha$, succeeds
  \item Return $\noe{6}(\text{base})$
\end{enumerate}

\textbf{Result:} $(\mathsf{Success}(\text{Feminine}(\text{base})), \{D6\})$
\end{example}

\section{Type Derivation Examples}

\begin{example}[Type Derivation for Fractal Application]
\label{ex:type-frac}
\textbf{Expression:} $e = \lambda f. \lambda x. f(\langle 1:4:7 \rangle(x))$

\textbf{Derivation Tree:}
\[
\infer[\textsc{T-Abs}]{\vdash \lambda f. \lambda x. f(\langle 1:4:7 \rangle(x)) : (\mathsf{Domain} \to \tau) \to \mathsf{Domain} \to \tau}{
  \infer[\textsc{T-Abs}]{f : \mathsf{Domain} \to \tau \vdash \lambda x. f(\langle 1:4:7 \rangle(x)) : \mathsf{Domain} \to \tau}{
    \infer[\textsc{T-App}]{f : \mathsf{Domain} \to \tau, x : \mathsf{Domain} \vdash f(\langle 1:4:7 \rangle(x)) : \tau}{
      \infer[\textsc{T-Var}]{f : \mathsf{Domain} \to \tau \vdash f : \mathsf{Domain} \to \tau}{}
      &
      \infer[\textsc{T-Frac}]{x : \mathsf{Domain} \vdash \langle 1:4:7 \rangle(x) : \mathsf{Domain}}{
        \infer[\textsc{T-Var}]{x : \mathsf{Domain} \vdash x : \mathsf{Domain}}{}
      }
    }
  }
}
\]

\textbf{Principal Type:} $\forall \tau. (\mathsf{Domain} \to \tau) \to \mathsf{Domain} \to \tau$
\end{example}

\begin{example}[Polymorphic RPM Function]
\label{ex:poly-rpm}
\textbf{Expression:}
\begin{verbatim}
let lift_with_check : forall a. Prereq -> a -> RPM[a] =
  /\a. \p : Prereq. \x : a.
    check(p) >>= \_ . return x
\end{verbatim}

\textbf{Type Derivation:}
\begin{align}
&\Gamma = \{a : \mathsf{Type}, p : \mathsf{Prereq}, x : a\} \\
&\Gamma \vdash \mathsf{check}(p) : \mathsf{RPM}[\mathsf{Unit}] \quad [\textsc{T-Check}] \\
&\Gamma \vdash \lambda\_. \rpmret\; x : \mathsf{Unit} \to \mathsf{RPM}[a] \quad [\textsc{T-Abs}, \textsc{T-Ret}] \\
&\Gamma \vdash \mathsf{check}(p) \rpmbind \lambda\_. \rpmret\; x : \mathsf{RPM}[a] \quad [\textsc{T-Bind}]
\end{align}

\textbf{Final Type:} $\forall a. \mathsf{Prereq} \to a \to \mathsf{RPM}[a]$
\end{example}

\section{Integrated Examples}

\begin{example}[Complete Healing Protocol]
\label{ex:healing-protocol}
\textbf{Goal:} Model a multi-phase healing transformation.

\textbf{Program:}
\begin{verbatim}
frac Diagnosis = <1:3>      -- Awareness + Below (examine)
frac Treatment = <4:8>      -- Vibration + Cause (apply)
frac Recovery = <7:9>       -- Rhythm + Effect (stabilize)

let healing_protocol : Domain -> RPM[Domain] =
  \patient : Domain .
    -- Phase 1: Diagnosis
    check(D2) >>= \_ .
    check(D3) >>= \_ .
    let diagnosed = frac Diagnosis (patient) in

    -- Phase 2: Treatment
    acquire(D2_treatment) >>= \_ .
    let treated = frac Treatment (diagnosed) in

    -- Phase 3: Recovery
    let recovered = frac Recovery (treated) in
    return recovered
\end{verbatim}

\textbf{Full Evaluation Trace:}

\textbf{Input:} $\text{patient} = D2_{3a}$ (Life, Mental, Negative---mental illness)
\textbf{Initial:} $\alpha_0 = \{D2, D3\}$

\textbf{Phase 1 - Diagnosis:}
\begin{align}
&\mathsf{check}(D2): D2 \in \alpha_0 \quad \checkmark \\
&\mathsf{check}(D3): D3 \in \alpha_0 \quad \checkmark \\
&\text{diagnosed} = \langle 1:3 \rangle(D2_{3a}) \\
&\quad = \noe{1}(\noe{3}(D2_{3a})) \\
&\quad = \noe{1}(\text{Below}(D2_{3a})) \\
&\quad = \text{Aware Below}(D2_{3a})
\end{align}

\textbf{Phase 2 - Treatment:}
\begin{align}
&\mathsf{acquire}(D2_{\text{treatment}}): \alpha_1 = \alpha_0 \cup \{D2_{\text{treatment}}\} \\
&\text{treated} = \langle 4:8 \rangle(\text{Aware Below } D2_{3a}) \\
&\quad = \noe{4}(\noe{8}(\text{Aware Below } D2_{3a})) \\
&\quad = \noe{4}(\text{Caused Aware Below } D2_{3a}) \\
&\quad = \text{Vibrating Caused Aware Below } D2_{3a}
\end{align}

\textbf{Phase 3 - Recovery:}
\begin{align}
&\text{recovered} = \langle 7:9 \rangle(\text{Vibrating Caused Aware Below } D2_{3a}) \\
&\quad = \noe{7}(\noe{9}(\text{Vibrating Caused Aware Below } D2_{3a})) \\
&\quad = \noe{7}(\text{Effected Vibrating Caused Aware Below } D2_{3a}) \\
&\quad = \text{Rhythmic Effected Vibrating Caused Aware Below } D2_{3a}
\end{align}

\textbf{Final Result:}
\[
(\mathsf{Success}(\text{Rhythmic Effected Vibrating Caused Aware Below } D2_{3a}), \alpha_1)
\]

\textbf{Interpretation:} The patient's mental illness has been:
\begin{itemize}
  \item Diagnosed (made conscious, grounded)
  \item Treated (energized, cause identified)
  \item Stabilized (rhythmic pattern, effect manifested)
\end{itemize}
\end{example}

\begin{example}[Wealth Manifestation with Prerequisites]
\label{ex:wealth-manifest}
\textbf{Goal:} Model wealth acquisition through the four worlds.

\textbf{Program:}
\begin{verbatim}
let manifest_wealth : RPM[Domain] =
  -- Spiritual foundation
  acquire(D1) >>= \_ .
  let spiritual = <1>(F6a) in

  -- Mental planning
  acquire(D3) >>= \_ .
  let mental = <1:4>(F6b) in

  -- Emotional alignment
  acquire(D4) >>= \_ .
  let emotional = <1:4:5:6>(F6c) in

  -- Physical manifestation
  check(D6) >>= \_ .
  let physical = <1:4:7:8:9>(F6d) in

  return physical
\end{verbatim}

\textbf{Evaluation (initial $\alpha = \{D6\}$):}
\begin{align}
&\mathsf{acquire}(D1): \alpha_1 = \{D6, D1\} \\
&\text{spiritual} = \noe{1}(F_{6a}) = \text{Aware Spiritual Wealth} \\[0.5em]
&\mathsf{acquire}(D3): \alpha_2 = \{D6, D1, D3\} \\
&\text{mental} = \langle 1:4 \rangle(F_{6b}) = \text{Aware Vibrating Mental Wealth} \\[0.5em]
&\mathsf{acquire}(D4): \alpha_3 = \{D6, D1, D3, D4\} \\
&\text{emotional} = \langle 1:4:5:6 \rangle(F_{6c}) = \text{Aware Vibrating Masc Fem Emotional Wealth} \\[0.5em]
&\mathsf{check}(D6): D6 \in \alpha_3 \quad \checkmark \\
&\text{physical} = \langle 1:4:7:8:9 \rangle(F_{6d}) \\
&\quad = \text{Aware Vibrating Rhythmic Caused Effected Physical Wealth}
\end{align}

\textbf{Final:} $(\mathsf{Success}(\text{Fully Transformed Physical Wealth}), \{D1, D3, D4, D6\})$
\end{example}

\begin{example}[Parallel Domain Development]
\label{ex:parallel-dev}
\textbf{Goal:} Develop Knowledge and Love domains concurrently.

\textbf{Program:}
\begin{verbatim}
let parallel_development : RPM[Domain * Domain] =
  let knowledge_path =
    acquire(D3) >>= \_ .
    return (<1:4:7>(D3_base))
  in
  let love_path =
    acquire(D4) >>= \_ .
    return (<1:4:7>(D4_base))
  in
  knowledge_path || love_path
\end{verbatim}

\textbf{Parallel Evaluation:}

Thread 1 (Knowledge):
\begin{align}
&\mathsf{acquire}(D3) \Rightarrow \alpha_1 = \{D3\} \\
&\langle 1:4:7 \rangle(D3_{\text{base}}) = \text{MVR}(D3_{\text{base}})
\end{align}

Thread 2 (Love):
\begin{align}
&\mathsf{acquire}(D4) \Rightarrow \alpha_2 = \{D4\} \\
&\langle 1:4:7 \rangle(D4_{\text{base}}) = \text{MVR}(D4_{\text{base}})
\end{align}

\textbf{Join:}
\[
\text{Result: } (\mathsf{Success}((\text{MVR}(D3), \text{MVR}(D4))), \{D3, D4\})
\]
\end{example}

\begin{example}[Quantum Superposition Example]
\label{ex:quantum-super}
\textbf{Goal:} Apply quantum MVR to a superposition state.

\textbf{Initial State:}
\[
\qket{\psi_0} = \frac{1}{\sqrt{2}}(\qket{A1} + \qket{D1})
\]
(Superposition of Spiritual and Physical Idea 1)

\textbf{Apply Quantum MVR:}
\begin{align}
\qop{\Frac}_{\text{MVR}} \qket{\psi_0} &= \qop{\nu}_1 \qop{\nu}_4 \qop{\nu}_7 \cdot \frac{1}{\sqrt{2}}(\qket{A1} + \qket{D1}) \\
&= \frac{1}{\sqrt{2}}(\qop{\Frac}_{\text{MVR}}\qket{A1} + \qop{\Frac}_{\text{MVR}}\qket{D1}) \\
&= \frac{1}{\sqrt{2}}(\qket{\text{MVR}(A1)} + \qket{\text{MVR}(D1)})
\end{align}

\textbf{Measurement:}
Observing ``which World?'' yields:
\begin{itemize}
  \item $P(\text{Spiritual}) = 1/2$
  \item $P(\text{Physical}) = 1/2$
\end{itemize}

\textbf{Interpretation:} The MVR transformation preserves the superposition structure until measurement.
\end{example}

\begin{example}[Denotational/Operational Equivalence]
\label{ex:den-op-equiv}
\textbf{Goal:} Verify semantic soundness for a specific expression.

\textbf{Expression:} $e = (\rpmret\; D1_{1a}) \rpmbind (\lambda x. \rpmret\; \noe{1}(x))$

\textbf{Operational (Big-Step):}
\begin{align}
&\langle (\rpmret\; D1_{1a}) \rpmbind (\lambda x. \rpmret\; \noe{1}(x)), \alpha \rangle \\
&\bigstep \langle (\lambda x. \rpmret\; \noe{1}(x))(D1_{1a}), \alpha \rangle \quad [\textsc{E-Ret-Bind}] \\
&\bigstep \langle \rpmret\; \noe{1}(D1_{1a}), \alpha \rangle \quad [\textsc{E-Beta}] \\
&\bigstep (\mathsf{Success}(\text{Aware } D1_{1a}), \alpha)
\end{align}

\textbf{Denotational:}
\begin{align}
\sem{e}_\rho(\alpha) &= \sem{(\rpmret\; D1_{1a}) \rpmbind (\lambda x. \rpmret\; \noe{1}(x))}_\rho(\alpha) \\
&= \sem{\lambda x. \rpmret\; \noe{1}(x)}_\rho(D1_{1a})(\alpha) \\
&= (\lambda v. \sem{\rpmret\; \noe{1}(v)}_\rho(\alpha))(D1_{1a}) \\
&= (\mathsf{Success}(\mathfrak{n}_1(D1_{1a})), \alpha) \\
&= (\mathsf{Success}(\text{Aware } D1_{1a}), \alpha)
\end{align}

\textbf{Verification:} Operational result = Denotational result $\checkmark$
\end{example}

\section{Additional Examples}

\begin{example}[Protection Protocol]
\label{ex:protection}
Apply protection fractal $\langle 5:7:8 \rangle$ (Masculine, Rhythm, Cause) to establish boundaries.
\end{example}

\begin{example}[Creative Expression]
\label{ex:creative}
Transform through D7 (Expression) using $\langle 1:5:6:4:9 \rangle$ for balanced creative output.
\end{example}

\begin{example}[Relationship Harmony]
\label{ex:relationship}
Use D4 (Love) with gender-balancing fractal $\langle 5:6:5:6 \rangle$ for partnership equilibrium.
\end{example}

\begin{example}[Learning Acceleration]
\label{ex:learning}
Apply D3 (Knowledge) fractal $\langle 1:1:4:4 \rangle$ for intensified awareness and absorption.
\end{example}

\begin{example}[Type Error Detection]
\label{ex:type-error}
Demonstrate how the type system catches invalid fractal indices or mismatched RPM types.
\end{example}

\chapter{Reference Tables}
\label{ch:reference}

\section{Domain Reference}

\begin{center}
\begin{longtable}{clll}
\caption{Complete Domain Reference} \\
\toprule
\textbf{ID} & \textbf{Domain} & \textbf{Foundations} & \textbf{Core Principle} \\
\midrule
\endfirsthead
\multicolumn{4}{c}{\textit{(continued)}} \\
\toprule
\textbf{ID} & \textbf{Domain} & \textbf{Foundations} & \textbf{Core Principle} \\
\midrule
\endhead
\midrule
\multicolumn{4}{r}{\textit{Continued on next page}} \\
\endfoot
\bottomrule
\endlastfoot
D1 & Spirit & Consciousness, Purpose, Connection & Awareness \\
D2 & Life & Health (Physical, Mental, Emotional) & Vitality \\
D3 & Knowledge & Learning, Wisdom, Understanding & Truth \\
D4 & Love & Relationships, Compassion, Unity & Connection \\
D5 & Power & Authority, Influence, Will & Mastery \\
D6 & Wealth & Resources, Abundance, Flow & Prosperity \\
D7 & Expression & Communication, Creativity, Art & Manifestation \\
\end{longtable}
\end{center}

\section{Noetic Function Reference}

\begin{center}
\begin{tabular}{cllll}
\toprule
\textbf{k} & \textbf{Name} & \textbf{Principle} & \textbf{Symbol} & \textbf{Action} \\
\midrule
0 & Identity & --- & $\id$ & No transformation \\
1 & Mentalism & Mind is All & $\mathfrak{n}_1$ & Apply awareness \\
2 & Correspondence (Above) & As above & $\mathfrak{n}_2$ & Project upward \\
3 & Correspondence (Below) & So below & $\mathfrak{n}_3$ & Ground downward \\
4 & Vibration & Everything moves & $\mathfrak{n}_4$ & Add energy/vibration \\
5 & Gender (Masculine) & Creation principle & $\mathfrak{n}_5$ & Active/projective \\
6 & Gender (Feminine) & Reception principle & $\mathfrak{n}_6$ & Receptive/formative \\
7 & Rhythm & Pendulum swings & $\mathfrak{n}_7$ & Establish pattern \\
8 & Cause & Every effect has cause & $\mathfrak{n}_8$ & Identify/create cause \\
9 & Effect & Every cause has effect & $\mathfrak{n}_9$ & Manifest effect \\
\bottomrule
\end{tabular}
\end{center}

\section{Standard Fractal Reference}

\begin{center}
\begin{tabular}{llll}
\toprule
\textbf{Name} & \textbf{Sequence} & \textbf{Class} & \textbf{Application} \\
\midrule
MVR & $\langle 1:4:7 \rangle$ & Stabilizing & General transformation \\
ACBE & $\langle 8:9 \rangle$ & Oscillatory & Causation analysis \\
Deep MVR & $\langle 1:1:1:4:4:4:7:7:7 \rangle$ & Stabilizing & Intensive work \\
Polarity & $\langle 2:3 \rangle$ & Oscillatory & Balance analysis \\
Gender & $\langle 5:6 \rangle$ & Oscillatory & Creative balance \\
Full Causation & $\langle 8:8:9:9 \rangle$ & Stabilizing & Deep causal analysis \\
Healing & $\langle 1:3:4:8:9 \rangle$ & Stabilizing & Health applications \\
Manifestation & $\langle 1:4:5:8:9:7 \rangle$ & Stabilizing & Creation/expression \\
\bottomrule
\end{tabular}
\end{center}

\section{Type Reference}

\begin{center}
\begin{tabular}{lll}
\toprule
\textbf{Type} & \textbf{Kind} & \textbf{Description} \\
\midrule
$\mathsf{Int}$ & $*$ & Integer values \\
$\mathsf{Bool}$ & $*$ & Boolean values \\
$\mathsf{Unit}$ & $*$ & Unit type (single value) \\
$\mathsf{Domain}$ & $*$ & Domain values \\
$\mathsf{Foundation}$ & $*$ & Foundation values \\
$\mathsf{Aspect}$ & $*$ & Aspect values \\
$\mathsf{Noe}[k]$ & $*$ & Noetic function type (indexed) \\
$\tau_1 \to \tau_2$ & $*$ & Function type \\
$\tau_1 \times \tau_2$ & $*$ & Product type \\
$\tau_1 + \tau_2$ & $*$ & Sum type \\
$\mathsf{RPM}[\tau]$ & $* \to *$ & RPM monad type \\
$\mathsf{Frac}[\bar{k}]$ & $*$ & Fractal type (indexed by sequence) \\
$\forall \alpha. \tau$ & $*$ & Polymorphic type \\
$\mathsf{List}[\tau]$ & $* \to *$ & List type \\
\bottomrule
\end{tabular}
\end{center}

%% ============================================================================
%% APPENDICES
%% ============================================================================
\appendix

\chapter{Symbol Index}
\label{app:symbols}

\begin{multicols}{2}
\begin{description}[leftmargin=2em,labelindent=0em]
\item[$\noe{k}$] Noetic function $k$
\item[$\mathfrak{n}_k$] Semantic value of $\noe{k}$
\item[$\Ideas$] Idea space
\item[$\Worlds$] World set
\item[$\Elements$] Element set
\item[$\Foundations$] Foundation set
\item[$\Acquisitions$] Acquisition set
\item[$\circ$] Function/fractal composition
\item[$\id$] Identity function/element
\item[$\RPM$] RPM monad type constructor
\item[$\rpmret$] Monadic return
\item[$\rpmbind$] Monadic bind
\item[$\Frac$] Fractal constructor
\item[$\FracSet$] Set of all fractals
\item[$\langle k_1 : \cdots : k_n \rangle$] Fractal notation
\item[$\Gamma$] Typing context
\item[$\vdash$] Typing judgment
\item[$\tau$] Type variable
\item[$\to$] Function type constructor
\item[$\times$] Product type constructor
\item[$+$] Sum type constructor
\item[$\sem{\cdot}$] Semantic brackets
\item[$\smallstep$] Small-step reduction
\item[$\bigstep$] Big-step evaluation
\item[$\AcqState$] Acquisition state
\item[$\Satisfied$] Prerequisite satisfied
\item[$\Failed$] Computation failed
\item[$\mathsf{check}$] Prerequisite check
\item[$\mathsf{acquire}$] Prerequisite acquisition
\item[$\eta$] Monad unit
\item[$\mu$] Monad multiplication
\item[$\text{ACBE}$] Functor $A \to D$
\item[$\Noetica$] Category Noetica
\item[$\Kleisli_\RPM$] RPM Kleisli category (v7.0)
\item[$\pcomp$] Parallel composition (v7.0)
\item[$\Hspace$] Hilbert space (v7.0)
\item[$\qket{\cdot}$] Quantum ket (v7.0)
\item[$\qbra{\cdot}$] Quantum bra (v7.0)
\item[$\subtype$] Subtype relation (v7.0)
\item[$\Spec$] Eigenvalue spectrum (v7.0)
\end{description}
\end{multicols}

\chapter{Complete Proof Details}
\label{app:proofs}

\section{Key Theorems and Their Proofs}

\begin{enumerate}
\item \textbf{Theorem (Monoid Structure):} $(\Noetics, \circ, \noe{0})$ forms a monoid. \textit{Proof:} Closure by composition of endofunctions; associativity by function composition; identity by $\noe{0}(x) = x$.

\item \textbf{Theorem (RPM Monad Laws):} $(\RPM, \eta, \mu)$ satisfies all monad laws. \textit{Proof:} See Chapter~\ref{ch:rpm-monad}, Theorem~\ref{thm:rpm-monad-laws}.

\item \textbf{Theorem (Fractal Composition Associativity):} $(\Frac_1 \circ \Frac_2) \circ \Frac_3 = \Frac_1 \circ (\Frac_2 \circ \Frac_3)$. \textit{Proof:} List concatenation is associative.

\item \textbf{Theorem (MVR Stability):} The MVR fractal is stabilizing. \textit{Proof:} Component analysis shows bounded iteration behavior.

\item \textbf{Theorem (Kleisli Category):} $\Kleisli_\RPM$ satisfies category axioms. \textit{Proof:} Follows from monad laws (Chapter~\ref{sec:rpm-kleisli}).

\item \textbf{Theorem (Progress):} Well-typed expressions are either values or can step. \textit{Proof:} Structural induction on typing derivations.

\item \textbf{Theorem (Preservation):} Reduction preserves types. \textit{Proof:} Case analysis on reduction rules with substitution lemma.

\item \textbf{Theorem (Type Safety):} Well-typed programs don't get stuck. \textit{Proof:} Corollary of Progress + Preservation.

\item \textbf{Theorem (Soundness):} Operational and denotational semantics agree. \textit{Proof:} Mutual induction on derivation structure.

\item \textbf{Theorem (Contraction Convergence):} Contracting fractals converge geometrically. \textit{Proof:} Application of Banach fixed-point theorem.
\end{enumerate}

\chapter{Implementation Notes}
\label{app:implementation}

\section{Interpreter Architecture}

A TKS interpreter should implement:

\begin{enumerate}
\item \textbf{Lexer/Parser:} According to grammar specifications
\item \textbf{Type Checker:} Implementing typing rules from the type system chapter
\item \textbf{Evaluator:} Following operational semantics
\item \textbf{RPM Runtime:} State monad with prerequisite tracking
\item \textbf{Fractal Engine:} Composition, simplification, and application
\end{enumerate}

\section{Suggested Implementation Order}

\begin{enumerate}
\item Core types and data structures
\item Noetic function implementations
\item Fractal composition and simplification
\item RPM monad with state threading
\item Type checker with inference
\item Big-step evaluator
\item Small-step reduction (for debugging)
\end{enumerate}

\section{Testing Strategy}

\begin{enumerate}
\item Unit tests for each noetic function
\item Property tests for monad laws
\item Fractal simplification tests
\item Integration tests for worked examples
\item Type safety tests (rejection of ill-typed programs)
\end{enumerate}

\chapter{Version History}
\label{app:versions}

\begin{description}
\item[v7.0 (Current)] Academic Expansion Release
  \begin{itemize}
  \item Complete Noetic Fractal Calculus with type theory and eigenfractals
  \item RPM Kleisli category and effect system
  \item Complete dynamic semantics for all constructs
  \item Full Progress and Preservation proofs
  \item Concurrency model prototype
  \item Quantum Noetic algebra prototype
  \item 20+ worked examples
  \end{itemize}

\item[v6.1] Dynamic Semantics Release
  \begin{itemize}
  \item RPM formalized as monad with proofs
  \item Basic Noetic Fractal Calculus
  \item Integrated evaluation rules
  \item Extended examples
  \end{itemize}

\item[v6.0] Academic Release

\item[v5.0] Mathematical Framework

\item[v3.2.1] Initial Formalization
\end{description}

%% ============================================================================
%% BACKMATTER
%% ============================================================================
\backmatter

\begin{thebibliography}{99}
\bibitem{kybalion} Three Initiates. \textit{The Kybalion}. Yogi Publication Society, 1908.
\bibitem{moggi} E. Moggi. ``Notions of computation and monads.'' \textit{Information and Computation}, 93(1):55--92, 1991.
\bibitem{wadler} P. Wadler. ``Monads for functional programming.'' \textit{Advanced Functional Programming}, LNCS 925:24--52, 1995.
\bibitem{pierce} B. C. Pierce. \textit{Types and Programming Languages}. MIT Press, 2002.
\bibitem{plotkin} G. D. Plotkin. ``A structural approach to operational semantics.'' DAIMI FN-19, 1981.
\bibitem{mac-lane} S. Mac Lane. \textit{Categories for the Working Mathematician}. Springer, 1971.
\bibitem{plotkin-power} G. Plotkin and J. Power. ``Algebraic operations and generic effects.'' \textit{Applied Categorical Structures}, 11(1):69--94, 2003.
\bibitem{nielsen-chuang} M. Nielsen and I. Chuang. \textit{Quantum Computation and Quantum Information}. Cambridge University Press, 2000.
\end{thebibliography}

\end{document}
